
\documentclass{article}
\usepackage{colortbl}
\usepackage{makecell}
\usepackage{multirow}
\usepackage{supertabular}

\begin{document}

\newcounter{utterance}

\centering \large Interaction Transcript for game `clean\_up', experiment `1\_medium\_7obj\_fr', episode 1 with nemotron{-}3{-}ultra.
\vspace{24pt}

{ \footnotesize  \setcounter{utterance}{1}
\setlength{\tabcolsep}{0pt}
\begin{supertabular}{c@{$\;$}|p{.15\linewidth}@{}p{.15\linewidth}p{.15\linewidth}p{.15\linewidth}p{.15\linewidth}p{.15\linewidth}}
    \# & $\;$A & \multicolumn{4}{c}{Game Master} & $\;\:$B\\
    \hline

    \theutterance \stepcounter{utterance}  
    & & \multicolumn{4}{p{0.6\linewidth}}{
        \cellcolor[rgb]{0.9,0.9,0.9}{
            \makecell[{{p{\linewidth}}}]{
                \texttt{\tiny{[P1$\langle$GM]}}
                \texttt{Je suis votre maître de jeu, et vous jouez à un jeu collaboratif avec la grille suivante comme plateau de jeu :} \\
\\ 
\texttt{```} \\
\texttt{ 1234567} \\
\texttt{╔══╤══╤═╗} \\
\texttt{║◌P│◌M│◌║ 1} \\
\texttt{╟──┼──┴─╢ 2} \\
\texttt{║◌◌│◌S◌O║ 3} \\
\texttt{╟──┘◌A◌◌║ 4} \\
\texttt{║D◌◌◌◌◌◌║ 5} \\
\texttt{║◌◌┌─┬──╢ 6} \\
\texttt{║T◌│◌│◌◌║ 7} \\
\texttt{╚══╧═╧══╝} \\
\\ 
\texttt{```} \\
\\ 
\texttt{* Le bord supérieur affiche des coordonnées x croissantes vers la droite, et le bord droit des coordonnées y croissantes vers le bas.} \\
\texttt{* Les objets suivants sont placés aléatoirement sur votre grille : 'P', 'O', 'T', 'S', 'D', 'A', 'M'.} \\
\\ 
\\ 
\texttt{L'autre joueur voit une variante du plateau de jeu, où les objets sont placés à des emplacements aléatoires différents. Vous ne pouvez pas voir le plateau de l'autre joueur et il ne peut pas voir le vôtre.} \\
\\ 
\\ 
\texttt{**Objectif :**} \\
\\ 
\texttt{Les deux joueurs doivent déplacer les objets sur leur arrière{-}plan respectif de manière à ce que des objets identiques se retrouvent aux mêmes coordonnées. Vous devez communiquer avec l'autre joueur pour convenir d'une configuration d'objectif commune.} \\
\\ 
\\ 
\texttt{**Règles :**} \\
\\ 
\texttt{* À chaque tour, vous pouvez envoyer exactement l'une des deux commandes suivantes :} \\
\texttt{ 1. `DIRE: <MESSAGE>` : pour envoyer un message (tout ce qui se trouve jusqu'au prochain saut de ligne) à l'autre joueur. Je le transmettrai à votre partenaire.} \\
\texttt{2. `DÉPLACER: <OBJET>, (<X>, <Y>)` : pour déplacer un objet vers une nouvelle position, où `<X>` est la colonne et `<Y>` est la ligne. Je vous informerai si votre déplacement a réussi ou non.} \\
\texttt{* Si vous ne respectez pas le format, ou si vous envoyez plusieurs commandes à la fois, je dois vous pénaliser.} \\
\texttt{* Si les deux joueurs accumulent plus de 16 pénalités, vous perdez tous les deux la partie.} \\
\texttt{* Il est essentiel que vous communiquiez avec l'autre joueur concernant votre état d'objectif ! La *seule* façon de transmettre votre stratégie à l'autre joueur est d'utiliser la commande `DIRE: <MESSAGE>` !} \\
\\ 
\\ 
\texttt{**Déplacement d'objets**} \\
\\ 
\texttt{* Vous ne pouvez déplacer des objets que vers des cellules situées à l'intérieur des limites de la grille. La cellule cible doit être vide, c'est{-}à{-}dire qu'elle ne doit contenir que le symbole '◌'.} \\
\texttt{* Si vous essayez de déplacer un objet vers un endroit qui n'est pas vide, ou si vous essayez de le déplacer en dehors de la grille, je dois vous pénaliser. Vous avez droit à un nouvel essai.} \\
\texttt{* Avant d'effectuer un déplacement, vérifiez bien que l'emplacement cible est vide et qu'il ne contient aucune lettre, aucun cadre ou aucune ligne !} \\
\\ 
\\ 
\texttt{**Fin de partie**} \\
\\ 
\texttt{Si vous pensez avoir atteint l'objectif d'aligner tous les objets, vous pouvez demander à l'autre joueur de terminer la partie en envoyant `DIRE: terminé ?`. Si l'autre joueur vous demande de terminer la partie et que vous répondez `DIRE: terminé !`, la partie se termine.} \\
\\ 
\texttt{Les deux joueurs gagnent si la partie se termine en 28 tours, un tour étant défini comme l'envoi d'une commande valide par chacun des deux joueurs.} \\
\\ 
\texttt{**Score :**} \\
\\ 
\texttt{Plus les objets identiques sont proches les uns des autres sur les deux plateaux de jeu, plus vous obtenez de points. Les pénalités réduisent vos points. Pouvez{-}vous battre le record ?} \\
\\ 
\\ 
\texttt{Envoyez un message à l'autre joueur pour commencer le jeu !} \\
            }
        }
    }
    & \\ \\

    \theutterance \stepcounter{utterance}  
    & \multicolumn{4}{p{0.6\linewidth}}{
        \cellcolor[rgb]{0.9,0.9,0.9}{
            \makecell[{{p{\linewidth}}}]{
                \texttt{\tiny{[P1$\rangle$GM]}}
                \texttt{DIRE: Bonjour ! Je vois les objets aux positions suivantes sur ma grille : A(3,4), D(1,5), M(3,1), O(7,3), P(2,1), S(4,3), T(1,7). Où sont les tiens ?} \\
            }
        }
    }
    & & \\ \\

    \theutterance \stepcounter{utterance}  
    & & & \multicolumn{2}{p{0.3\linewidth}}{
        \cellcolor[rgb]{0.9,0.9,0.9}{
            \makecell[{{p{\linewidth}}}]{
                \texttt{\tiny{[GM$|$GM]}}
                \texttt{DIRE: Bonjour ! Je vois les objets aux positions suivantes sur ma grille : A(3,4), D(1,5), M(3,1), O(7,3), P(2,1), S(4,3), T(1,7). Où sont les tiens ?} \\
            }
        }
    }
    & & \\ \\

    \theutterance \stepcounter{utterance}  
    & & \multicolumn{4}{p{0.6\linewidth}}{
        \cellcolor[rgb]{0.9,0.9,0.9}{
            \makecell[{{p{\linewidth}}}]{
                \texttt{\tiny{[GM$\rangle$P2]}}
                \texttt{Je suis votre maître de jeu, et vous jouez à un jeu collaboratif avec la grille suivante comme plateau de jeu :} \\
\\ 
\texttt{```} \\
\texttt{ 1234567} \\
\texttt{╔══╤══╤═╗} \\
\texttt{║◌◌│◌◌│T║ 1} \\
\texttt{╟──┼──┴─╢ 2} \\
\texttt{║◌◌│◌S◌P║ 3} \\
\texttt{╟──┘◌◌MO║ 4} \\
\texttt{║D◌◌◌◌◌◌║ 5} \\
\texttt{║◌◌┌─┬──╢ 6} \\
\texttt{║◌◌│A│◌◌║ 7} \\
\texttt{╚══╧═╧══╝} \\
\\ 
\texttt{```} \\
\\ 
\texttt{* Le bord supérieur affiche des coordonnées x croissantes vers la droite, et le bord droit des coordonnées y croissantes vers le bas.} \\
\texttt{* Les objets suivants sont placés aléatoirement sur votre grille : 'P', 'O', 'T', 'S', 'D', 'A', 'M'.} \\
\\ 
\\ 
\texttt{L'autre joueur voit une variante du plateau de jeu, où les objets sont placés à des emplacements aléatoires différents. Vous ne pouvez pas voir le plateau de l'autre joueur et il ne peut pas voir le vôtre.} \\
\\ 
\\ 
\texttt{**Objectif :**} \\
\\ 
\texttt{Les deux joueurs doivent déplacer les objets sur leur arrière{-}plan respectif de manière à ce que des objets identiques se retrouvent aux mêmes coordonnées. Vous devez communiquer avec l'autre joueur pour convenir d'une configuration d'objectif commune.} \\
\\ 
\\ 
\texttt{**Règles :**} \\
\\ 
\texttt{* À chaque tour, vous pouvez envoyer exactement l'une des deux commandes suivantes :} \\
\texttt{ 1. `DIRE: <MESSAGE>` : pour envoyer un message (tout ce qui se trouve jusqu'au prochain saut de ligne) à l'autre joueur. Je le transmettrai à votre partenaire.} \\
\texttt{2. `DÉPLACER: <OBJET>, (<X>, <Y>)` : pour déplacer un objet vers une nouvelle position, où `<X>` est la colonne et `<Y>` est la ligne. Je vous informerai si votre déplacement a réussi ou non.} \\
\texttt{* Si vous ne respectez pas le format, ou si vous envoyez plusieurs commandes à la fois, je dois vous pénaliser.} \\
\texttt{* Si les deux joueurs accumulent plus de 16 pénalités, vous perdez tous les deux la partie.} \\
\texttt{* Il est essentiel que vous communiquiez avec l'autre joueur concernant votre état d'objectif ! La *seule* façon de transmettre votre stratégie à l'autre joueur est d'utiliser la commande `DIRE: <MESSAGE>` !} \\
\\ 
\\ 
\texttt{**Déplacement d'objets**} \\
\\ 
\texttt{* Vous ne pouvez déplacer des objets que vers des cellules situées à l'intérieur des limites de la grille. La cellule cible doit être vide, c'est{-}à{-}dire qu'elle ne doit contenir que le symbole '◌'.} \\
\texttt{* Si vous essayez de déplacer un objet vers un endroit qui n'est pas vide, ou si vous essayez de le déplacer en dehors de la grille, je dois vous pénaliser. Vous avez droit à un nouvel essai.} \\
\texttt{* Avant d'effectuer un déplacement, vérifiez bien que l'emplacement cible est vide et qu'il ne contient aucune lettre, aucun cadre ou aucune ligne !} \\
\\ 
\\ 
\texttt{**Fin de partie**} \\
\\ 
\texttt{Si vous pensez avoir atteint l'objectif d'aligner tous les objets, vous pouvez demander à l'autre joueur de terminer la partie en envoyant `DIRE: terminé ?`. Si l'autre joueur vous demande de terminer la partie et que vous répondez `DIRE: terminé !`, la partie se termine.} \\
\\ 
\texttt{Les deux joueurs gagnent si la partie se termine en 28 tours, un tour étant défini comme l'envoi d'une commande valide par chacun des deux joueurs.} \\
\\ 
\texttt{**Score :**} \\
\\ 
\texttt{Plus les objets identiques sont proches les uns des autres sur les deux plateaux de jeu, plus vous obtenez de points. Les pénalités réduisent vos points. Pouvez{-}vous battre le record ?} \\
\\ 
\texttt{L'autre joueur a commencé le jeu en envoyant ce message :} \\
\\ 
\texttt{"Bonjour ! Je vois les objets aux positions suivantes sur ma grille : A(3,4), D(1,5), M(3,1), O(7,3), P(2,1), S(4,3), T(1,7). Où sont les tiens ?"} \\
\\ 
\texttt{Quelle est votre première commande ?} \\
            }
        }
    }
    & \\ \\

    \theutterance \stepcounter{utterance}  
    & & & \multicolumn{4}{p{0.6\linewidth}}{
        \cellcolor[rgb]{0.9,0.9,0.9}{
            \makecell[{{p{\linewidth}}}]{
                \texttt{\tiny{[GM$\langle$P2]}}
                \texttt{DIRE: Bonjour ! Je vois les objets aux positions suivantes sur ma grille : T(3,1), S(4,3), P(7,3), M(6,4), O(7,4), D(1,5), A(3,7).} \\
\\ 
\texttt{On doit se mettre d'accord sur une configuration cible commune. Pour faciliter les déplacements, je propose de rassembler tous les objets dans le coin bas{-}gauche (colonnes 1{-}2, lignes 5{-}7) qui semble dégagé chez moi. Par exemple :} \\
\texttt{{-} D(1,5) (déjà en place)} \\
\texttt{{-} A(1,6)} \\
\texttt{{-} M(2,5)} \\
\texttt{{-} O(2,6)} \\
\texttt{{-} P(1,7)} \\
\texttt{{-} S(2,7)} \\
\texttt{{-} T(2,6) ? (conflit avec O)} \\
\\ 
\texttt{Qu'en penses{-}tu ? As{-}tu de la place dans cette zone chez toi ?} \\
            }
        }
    }
     \\ \\

    \theutterance \stepcounter{utterance}  
    & & & \multicolumn{2}{p{0.3\linewidth}}{
        \cellcolor[rgb]{0.9,0.9,0.9}{
            \makecell[{{p{\linewidth}}}]{
                \texttt{\tiny{[GM$|$GM]}}
                \texttt{DIRE: Bonjour ! Je vois les objets aux positions suivantes sur ma grille : T(3,1), S(4,3), P(7,3), M(6,4), O(7,4), D(1,5), A(3,7).} \\
\\ 
\texttt{On doit se mettre d'accord sur une configuration cible commune. Pour faciliter les déplacements, je propose de rassembler tous les objets dans le coin bas{-}gauche (colonnes 1{-}2, lignes 5{-}7) qui semble dégagé chez moi. Par exemple :} \\
\texttt{{-} D(1,5) (déjà en place)} \\
\texttt{{-} A(1,6)} \\
\texttt{{-} M(2,5)} \\
\texttt{{-} O(2,6)} \\
\texttt{{-} P(1,7)} \\
\texttt{{-} S(2,7)} \\
\texttt{{-} T(2,6) ? (conflit avec O)} \\
\\ 
\texttt{Qu'en penses{-}tu ? As{-}tu de la place dans cette zone chez toi ?} \\
            }
        }
    }
    & & \\ \\

    \theutterance \stepcounter{utterance}  
    & & & \multicolumn{2}{p{0.3\linewidth}}{
        \cellcolor[rgb]{0.9,0.9,0.9}{
            \makecell[{{p{\linewidth}}}]{
                \texttt{\tiny{[GM$|$GM]}}
                \texttt{Invalid format: tail is not empty: '} \\
\\ 
\texttt{On doit se mettre d'accord sur une configuration cible commune. Pour faciliter les déplacements, je propose de rassembler tous les objets dans le coin bas{-}gauche (colonnes 1{-}2, lignes 5{-}7) qui semble dégagé chez moi. Par exemple :} \\
\texttt{{-} D(1,5) (déjà en place)} \\
\texttt{{-} A(1,6)} \\
\texttt{{-} M(2,5)} \\
\texttt{{-} O(2,6)} \\
\texttt{{-} P(1,7)} \\
\texttt{{-} S(2,7)} \\
\texttt{{-} T(2,6) ? (conflit avec O)} \\
\\ 
\texttt{Qu'en penses{-}tu ? As{-}tu de la place dans cette zone chez toi ?} \\
            }
        }
    }
    & & \\ \\

    \theutterance \stepcounter{utterance}  
    & & \multicolumn{4}{p{0.6\linewidth}}{
        \cellcolor[rgb]{0.9,0.9,0.9}{
            \makecell[{{p{\linewidth}}}]{
                \texttt{\tiny{[GM$\rangle$P2]}}
                \texttt{Pénalité : Votre message ne doit rien contenir après la commande !} \\
\\ 
\texttt{Assurez{-}vous que votre réponse ne contienne que `DIRE: <MESSAGE>` ou `DÉPLACER: <OBJET>, (<X>, <Y>)`, et rien d'autre !} \\
\\ 
\texttt{Vous avez collectivement accumulé 1 pénalités sur 16.} \\
\texttt{Veuillez réessayer !} \\
            }
        }
    }
    & \\ \\

    \theutterance \stepcounter{utterance}  
    & & & \multicolumn{4}{p{0.6\linewidth}}{
        \cellcolor[rgb]{0.9,0.9,0.9}{
            \makecell[{{p{\linewidth}}}]{
                \texttt{\tiny{[GM$\langle$P2]}}
                \texttt{DIRE: Bonjour ! Je vois les objets aux positions suivantes sur ma grille : T(3,1), S(4,3), P(7,3), M(6,4), O(7,4), D(1,5), A(3,7). On doit se mettre d'accord sur une configuration cible commune. Pour faciliter les déplacements, je propose de rassembler tous les objets dans le coin bas{-}gauche (colonnes 1{-}2, lignes 5{-}7) qui semble dégagé chez moi. Par exemple : D(1,5), A(1,6), M(2,5), O(2,6), P(1,7), S(2,7), T(1,6) ? Qu'en penses{-}tu ? As{-}tu de la place dans cette zone chez toi ?} \\
            }
        }
    }
     \\ \\

    \theutterance \stepcounter{utterance}  
    & & & \multicolumn{2}{p{0.3\linewidth}}{
        \cellcolor[rgb]{0.9,0.9,0.9}{
            \makecell[{{p{\linewidth}}}]{
                \texttt{\tiny{[GM$|$GM]}}
                \texttt{DIRE: Bonjour ! Je vois les objets aux positions suivantes sur ma grille : T(3,1), S(4,3), P(7,3), M(6,4), O(7,4), D(1,5), A(3,7). On doit se mettre d'accord sur une configuration cible commune. Pour faciliter les déplacements, je propose de rassembler tous les objets dans le coin bas{-}gauche (colonnes 1{-}2, lignes 5{-}7) qui semble dégagé chez moi. Par exemple : D(1,5), A(1,6), M(2,5), O(2,6), P(1,7), S(2,7), T(1,6) ? Qu'en penses{-}tu ? As{-}tu de la place dans cette zone chez toi ?} \\
            }
        }
    }
    & & \\ \\

    \theutterance \stepcounter{utterance}  
    & & \multicolumn{4}{p{0.6\linewidth}}{
        \cellcolor[rgb]{0.9,0.9,0.9}{
            \makecell[{{p{\linewidth}}}]{
                \texttt{\tiny{[P1$\langle$GM]}}
                \texttt{Votre message a été transmis à l'autre joueur.} \\
\\ 
\\ 
\texttt{Vous jouez actuellement le tour 1 sur un maximum de 28.} \\
\texttt{Vous avez collectivement accumulé 1 pénalités sur 16.} \\
\texttt{L'autre joueur a envoyé ce message :} \\
\\ 
\\ 
\texttt{"Bonjour ! Je vois les objets aux positions suivantes sur ma grille : T(3,1), S(4,3), P(7,3), M(6,4), O(7,4), D(1,5), A(3,7). On doit se mettre d'accord sur une configuration cible commune. Pour faciliter les déplacements, je propose de rassembler tous les objets dans le coin bas{-}gauche (colonnes 1{-}2, lignes 5{-}7) qui semble dégagé chez moi. Par exemple : D(1,5), A(1,6), M(2,5), O(2,6), P(1,7), S(2,7), T(1,6) ? Qu'en penses{-}tu ? As{-}tu de la place dans cette zone chez toi ?"} \\
\\ 
\\ 
\texttt{Quelle est votre prochaine commande ?} \\
            }
        }
    }
    & \\ \\

    \theutterance \stepcounter{utterance}  
    & \multicolumn{4}{p{0.6\linewidth}}{
        \cellcolor[rgb]{0.9,0.9,0.9}{
            \makecell[{{p{\linewidth}}}]{
                \texttt{\tiny{[P1$\rangle$GM]}}
                \texttt{DIRE: Chez moi, la zone bas{-}gauche (colonnes 1{-}2, lignes 5{-}7) n'est pas libre : D est déjà en (1,5), T en (1,7), et A en (3,4) bloque l'accès. En revanche, le coin bas{-}droit (colonnes 6{-}7, lignes 5{-}7) a l'air vide (juste T en (1,7) loin de là). As{-}tu de la place là{-}bas chez toi ? Sinon, on peut viser le coin haut{-}droit (colonnes 6{-}7, lignes 1{-}3) qui est libre chez moi. Quelles sont tes cases vides dans ces coins ?} \\
            }
        }
    }
    & & \\ \\

    \theutterance \stepcounter{utterance}  
    & & & \multicolumn{2}{p{0.3\linewidth}}{
        \cellcolor[rgb]{0.9,0.9,0.9}{
            \makecell[{{p{\linewidth}}}]{
                \texttt{\tiny{[GM$|$GM]}}
                \texttt{DIRE: Chez moi, la zone bas{-}gauche (colonnes 1{-}2, lignes 5{-}7) n'est pas libre : D est déjà en (1,5), T en (1,7), et A en (3,4) bloque l'accès. En revanche, le coin bas{-}droit (colonnes 6{-}7, lignes 5{-}7) a l'air vide (juste T en (1,7) loin de là). As{-}tu de la place là{-}bas chez toi ? Sinon, on peut viser le coin haut{-}droit (colonnes 6{-}7, lignes 1{-}3) qui est libre chez moi. Quelles sont tes cases vides dans ces coins ?} \\
            }
        }
    }
    & & \\ \\

    \theutterance \stepcounter{utterance}  
    & & \multicolumn{4}{p{0.6\linewidth}}{
        \cellcolor[rgb]{0.9,0.9,0.9}{
            \makecell[{{p{\linewidth}}}]{
                \texttt{\tiny{[GM$\rangle$P2]}}
                \texttt{Votre message a été transmis à l'autre joueur.} \\
\\ 
\\ 
\texttt{Vous jouez actuellement le tour 2 sur un maximum de 28.} \\
\texttt{Vous avez collectivement accumulé 1 pénalités sur 16.} \\
\texttt{L'autre joueur a envoyé ce message :} \\
\\ 
\\ 
\texttt{"Chez moi, la zone bas{-}gauche (colonnes 1{-}2, lignes 5{-}7) n'est pas libre : D est déjà en (1,5), T en (1,7), et A en (3,4) bloque l'accès. En revanche, le coin bas{-}droit (colonnes 6{-}7, lignes 5{-}7) a l'air vide (juste T en (1,7) loin de là). As{-}tu de la place là{-}bas chez toi ? Sinon, on peut viser le coin haut{-}droit (colonnes 6{-}7, lignes 1{-}3) qui est libre chez moi. Quelles sont tes cases vides dans ces coins ?"} \\
\\ 
\\ 
\texttt{Quelle est votre prochaine commande ?} \\
            }
        }
    }
    & \\ \\

    \theutterance \stepcounter{utterance}  
    & & & \multicolumn{4}{p{0.6\linewidth}}{
        \cellcolor[rgb]{0.9,0.9,0.9}{
            \makecell[{{p{\linewidth}}}]{
                \texttt{\tiny{[GM$\langle$P2]}}
                \texttt{DIRE: Chez moi, le coin bas{-}droit (col 6{-}7, lig 5{-}7) a M(6,4) et O(7,4) juste au{-}dessus, mais les cases (6,5),(7,5),(6,6),(7,6),(6,7),(7,7) semblent vides. Le coin haut{-}droit (col 6{-}7, lig 1{-}3) a P(7,3) et T(3,1) loin, donc (6,1),(7,1),(6,2),(7,2),(6,3) sont libres. Le coin haut{-}droit a l'air plus facile : 5 cases libres contiguës. On vise là{-}bas ? Configuration proposée : T(6,1), P(7,1), M(6,2), O(7,2), S(6,3), A(7,3), D(6,4) ? Dis{-}moi si ça te va.} \\
            }
        }
    }
     \\ \\

    \theutterance \stepcounter{utterance}  
    & & & \multicolumn{2}{p{0.3\linewidth}}{
        \cellcolor[rgb]{0.9,0.9,0.9}{
            \makecell[{{p{\linewidth}}}]{
                \texttt{\tiny{[GM$|$GM]}}
                \texttt{DIRE: Chez moi, le coin bas{-}droit (col 6{-}7, lig 5{-}7) a M(6,4) et O(7,4) juste au{-}dessus, mais les cases (6,5),(7,5),(6,6),(7,6),(6,7),(7,7) semblent vides. Le coin haut{-}droit (col 6{-}7, lig 1{-}3) a P(7,3) et T(3,1) loin, donc (6,1),(7,1),(6,2),(7,2),(6,3) sont libres. Le coin haut{-}droit a l'air plus facile : 5 cases libres contiguës. On vise là{-}bas ? Configuration proposée : T(6,1), P(7,1), M(6,2), O(7,2), S(6,3), A(7,3), D(6,4) ? Dis{-}moi si ça te va.} \\
            }
        }
    }
    & & \\ \\

    \theutterance \stepcounter{utterance}  
    & & \multicolumn{4}{p{0.6\linewidth}}{
        \cellcolor[rgb]{0.9,0.9,0.9}{
            \makecell[{{p{\linewidth}}}]{
                \texttt{\tiny{[P1$\langle$GM]}}
                \texttt{Votre message a été transmis à l'autre joueur.} \\
\\ 
\\ 
\texttt{Vous jouez actuellement le tour 2 sur un maximum de 28.} \\
\texttt{Vous avez collectivement accumulé 1 pénalités sur 16.} \\
\texttt{L'autre joueur a envoyé ce message :} \\
\\ 
\\ 
\texttt{"Chez moi, le coin bas{-}droit (col 6{-}7, lig 5{-}7) a M(6,4) et O(7,4) juste au{-}dessus, mais les cases (6,5),(7,5),(6,6),(7,6),(6,7),(7,7) semblent vides. Le coin haut{-}droit (col 6{-}7, lig 1{-}3) a P(7,3) et T(3,1) loin, donc (6,1),(7,1),(6,2),(7,2),(6,3) sont libres. Le coin haut{-}droit a l'air plus facile : 5 cases libres contiguës. On vise là{-}bas ? Configuration proposée : T(6,1), P(7,1), M(6,2), O(7,2), S(6,3), A(7,3), D(6,4) ? Dis{-}moi si ça te va."} \\
\\ 
\\ 
\texttt{Quelle est votre prochaine commande ?} \\
            }
        }
    }
    & \\ \\

    \theutterance \stepcounter{utterance}  
    & \multicolumn{4}{p{0.6\linewidth}}{
        \cellcolor[rgb]{0.9,0.9,0.9}{
            \makecell[{{p{\linewidth}}}]{
                \texttt{\tiny{[P1$\rangle$GM]}}
                \texttt{DIRE: Le coin haut{-}droit (colonnes 6{-}7, lignes 1{-}4) est parfait chez moi, tout est vide. Ta configuration proposée me va : T(6,1), P(7,1), M(6,2), O(7,2), S(6,3), A(7,3), D(6,4). On y va ! Je commence par déplacer T vers (6,1).} \\
            }
        }
    }
    & & \\ \\

    \theutterance \stepcounter{utterance}  
    & & & \multicolumn{2}{p{0.3\linewidth}}{
        \cellcolor[rgb]{0.9,0.9,0.9}{
            \makecell[{{p{\linewidth}}}]{
                \texttt{\tiny{[GM$|$GM]}}
                \texttt{DIRE: Le coin haut{-}droit (colonnes 6{-}7, lignes 1{-}4) est parfait chez moi, tout est vide. Ta configuration proposée me va : T(6,1), P(7,1), M(6,2), O(7,2), S(6,3), A(7,3), D(6,4). On y va ! Je commence par déplacer T vers (6,1).} \\
            }
        }
    }
    & & \\ \\

    \theutterance \stepcounter{utterance}  
    & & \multicolumn{4}{p{0.6\linewidth}}{
        \cellcolor[rgb]{0.9,0.9,0.9}{
            \makecell[{{p{\linewidth}}}]{
                \texttt{\tiny{[GM$\rangle$P2]}}
                \texttt{Votre message a été transmis à l'autre joueur.} \\
\\ 
\\ 
\texttt{Vous jouez actuellement le tour 3 sur un maximum de 28.} \\
\texttt{Vous avez collectivement accumulé 1 pénalités sur 16.} \\
\texttt{L'autre joueur a envoyé ce message :} \\
\\ 
\\ 
\texttt{"Le coin haut{-}droit (colonnes 6{-}7, lignes 1{-}4) est parfait chez moi, tout est vide. Ta configuration proposée me va : T(6,1), P(7,1), M(6,2), O(7,2), S(6,3), A(7,3), D(6,4). On y va ! Je commence par déplacer T vers (6,1)."} \\
\\ 
\\ 
\texttt{Quelle est votre prochaine commande ?} \\
            }
        }
    }
    & \\ \\

    \theutterance \stepcounter{utterance}  
    & & & \multicolumn{4}{p{0.6\linewidth}}{
        \cellcolor[rgb]{0.9,0.9,0.9}{
            \makecell[{{p{\linewidth}}}]{
                \texttt{\tiny{[GM$\langle$P2]}}
                \texttt{DÉPLACER: T, (6, 1)} \\
            }
        }
    }
     \\ \\

    \theutterance \stepcounter{utterance}  
    & & & \multicolumn{2}{p{0.3\linewidth}}{
        \cellcolor[rgb]{0.9,0.9,0.9}{
            \makecell[{{p{\linewidth}}}]{
                \texttt{\tiny{[GM$|$GM]}}
                \texttt{DÉPLACER: T, (6, 1)} \\
            }
        }
    }
    & & \\ \\

    \theutterance \stepcounter{utterance}  
    & & & \multicolumn{2}{p{0.3\linewidth}}{
        \cellcolor[rgb]{0.9,0.9,0.9}{
            \makecell[{{p{\linewidth}}}]{
                \texttt{\tiny{[GM$|$GM]}}
                \texttt{Pénalité : (6,1) n'est pas vide, mais contient '│'.} \\
\\ 
\texttt{Vous avez collectivement accumulé 2 pénalités sur 16.} \\
\texttt{Veuillez réessayer !} \\
            }
        }
    }
    & & \\ \\

    \theutterance \stepcounter{utterance}  
    & & \multicolumn{4}{p{0.6\linewidth}}{
        \cellcolor[rgb]{0.9,0.9,0.9}{
            \makecell[{{p{\linewidth}}}]{
                \texttt{\tiny{[GM$\rangle$P2]}}
                \texttt{Pénalité : (6,1) n'est pas vide, mais contient '│'.} \\
\\ 
\texttt{Vous avez collectivement accumulé 2 pénalités sur 16.} \\
\texttt{Veuillez réessayer !} \\
            }
        }
    }
    & \\ \\

    \theutterance \stepcounter{utterance}  
    & & & \multicolumn{4}{p{0.6\linewidth}}{
        \cellcolor[rgb]{0.9,0.9,0.9}{
            \makecell[{{p{\linewidth}}}]{
                \texttt{\tiny{[GM$\langle$P2]}}
                \texttt{DÉPLACER: T, (7, 1)} \\
            }
        }
    }
     \\ \\

    \theutterance \stepcounter{utterance}  
    & & & \multicolumn{2}{p{0.3\linewidth}}{
        \cellcolor[rgb]{0.9,0.9,0.9}{
            \makecell[{{p{\linewidth}}}]{
                \texttt{\tiny{[GM$|$GM]}}
                \texttt{DÉPLACER: T, (7, 1)} \\
            }
        }
    }
    & & \\ \\

    \theutterance \stepcounter{utterance}  
    & & & \multicolumn{2}{p{0.3\linewidth}}{
        \cellcolor[rgb]{0.9,0.9,0.9}{
            \makecell[{{p{\linewidth}}}]{
                \texttt{\tiny{[GM$|$GM]}}
                \texttt{Pénalité : (7,1) n'est pas vide, mais contient 'T'.} \\
\\ 
\texttt{Vous avez collectivement accumulé 3 pénalités sur 16.} \\
\texttt{Veuillez réessayer !} \\
            }
        }
    }
    & & \\ \\

    \theutterance \stepcounter{utterance}  
    & & \multicolumn{4}{p{0.6\linewidth}}{
        \cellcolor[rgb]{0.9,0.9,0.9}{
            \makecell[{{p{\linewidth}}}]{
                \texttt{\tiny{[GM$\rangle$P2]}}
                \texttt{Pénalité : (7,1) n'est pas vide, mais contient 'T'.} \\
\\ 
\texttt{Vous avez collectivement accumulé 3 pénalités sur 16.} \\
\texttt{Veuillez réessayer !} \\
            }
        }
    }
    & \\ \\

    \theutterance \stepcounter{utterance}  
    & & & \multicolumn{4}{p{0.6\linewidth}}{
        \cellcolor[rgb]{0.9,0.9,0.9}{
            \makecell[{{p{\linewidth}}}]{
                \texttt{\tiny{[GM$\langle$P2]}}
                \texttt{DÉPLACER: T, (6, 2)} \\
            }
        }
    }
     \\ \\

    \theutterance \stepcounter{utterance}  
    & & & \multicolumn{2}{p{0.3\linewidth}}{
        \cellcolor[rgb]{0.9,0.9,0.9}{
            \makecell[{{p{\linewidth}}}]{
                \texttt{\tiny{[GM$|$GM]}}
                \texttt{DÉPLACER: T, (6, 2)} \\
            }
        }
    }
    & & \\ \\

    \theutterance \stepcounter{utterance}  
    & & & \multicolumn{2}{p{0.3\linewidth}}{
        \cellcolor[rgb]{0.9,0.9,0.9}{
            \makecell[{{p{\linewidth}}}]{
                \texttt{\tiny{[GM$|$GM]}}
                \texttt{Pénalité : (6,2) n'est pas vide, mais contient '┴'.} \\
\\ 
\texttt{Vous avez collectivement accumulé 4 pénalités sur 16.} \\
\texttt{Veuillez réessayer !} \\
            }
        }
    }
    & & \\ \\

    \theutterance \stepcounter{utterance}  
    & & \multicolumn{4}{p{0.6\linewidth}}{
        \cellcolor[rgb]{0.9,0.9,0.9}{
            \makecell[{{p{\linewidth}}}]{
                \texttt{\tiny{[GM$\rangle$P2]}}
                \texttt{Pénalité : (6,2) n'est pas vide, mais contient '┴'.} \\
\\ 
\texttt{Vous avez collectivement accumulé 4 pénalités sur 16.} \\
\texttt{Veuillez réessayer !} \\
            }
        }
    }
    & \\ \\

    \theutterance \stepcounter{utterance}  
    & & & \multicolumn{4}{p{0.6\linewidth}}{
        \cellcolor[rgb]{0.9,0.9,0.9}{
            \makecell[{{p{\linewidth}}}]{
                \texttt{\tiny{[GM$\langle$P2]}}
                \texttt{DÉPLACER: T, (7, 2)} \\
            }
        }
    }
     \\ \\

    \theutterance \stepcounter{utterance}  
    & & & \multicolumn{2}{p{0.3\linewidth}}{
        \cellcolor[rgb]{0.9,0.9,0.9}{
            \makecell[{{p{\linewidth}}}]{
                \texttt{\tiny{[GM$|$GM]}}
                \texttt{DÉPLACER: T, (7, 2)} \\
            }
        }
    }
    & & \\ \\

    \theutterance \stepcounter{utterance}  
    & & & \multicolumn{2}{p{0.3\linewidth}}{
        \cellcolor[rgb]{0.9,0.9,0.9}{
            \makecell[{{p{\linewidth}}}]{
                \texttt{\tiny{[GM$|$GM]}}
                \texttt{Pénalité : (7,2) n'est pas vide, mais contient '─'.} \\
\\ 
\texttt{Vous avez collectivement accumulé 5 pénalités sur 16.} \\
\texttt{Veuillez réessayer !} \\
            }
        }
    }
    & & \\ \\

    \theutterance \stepcounter{utterance}  
    & & \multicolumn{4}{p{0.6\linewidth}}{
        \cellcolor[rgb]{0.9,0.9,0.9}{
            \makecell[{{p{\linewidth}}}]{
                \texttt{\tiny{[GM$\rangle$P2]}}
                \texttt{Pénalité : (7,2) n'est pas vide, mais contient '─'.} \\
\\ 
\texttt{Vous avez collectivement accumulé 5 pénalités sur 16.} \\
\texttt{Veuillez réessayer !} \\
            }
        }
    }
    & \\ \\

    \theutterance \stepcounter{utterance}  
    & & & \multicolumn{4}{p{0.6\linewidth}}{
        \cellcolor[rgb]{0.9,0.9,0.9}{
            \makecell[{{p{\linewidth}}}]{
                \texttt{\tiny{[GM$\langle$P2]}}
                \texttt{DÉPLACER: T, (6, 3)} \\
            }
        }
    }
     \\ \\

    \theutterance \stepcounter{utterance}  
    & & & \multicolumn{2}{p{0.3\linewidth}}{
        \cellcolor[rgb]{0.9,0.9,0.9}{
            \makecell[{{p{\linewidth}}}]{
                \texttt{\tiny{[GM$|$GM]}}
                \texttt{DÉPLACER: T, (6, 3)} \\
            }
        }
    }
    & & \\ \\

    \theutterance \stepcounter{utterance}  
    & & & \multicolumn{2}{p{0.3\linewidth}}{
        \cellcolor[rgb]{0.9,0.9,0.9}{
            \makecell[{{p{\linewidth}}}]{
                \texttt{\tiny{[GM$|$GM]}}
                \texttt{T moved to (6, 3)} \\
            }
        }
    }
    & & \\ \\

    \theutterance \stepcounter{utterance}  
    & & \multicolumn{4}{p{0.6\linewidth}}{
        \cellcolor[rgb]{0.9,0.9,0.9}{
            \makecell[{{p{\linewidth}}}]{
                \texttt{\tiny{[P1$\langle$GM]}}
                \texttt{Votre message a été transmis à l'autre joueur.} \\
\\ 
\\ 
\texttt{Vous jouez actuellement le tour 3 sur un maximum de 28.} \\
\texttt{Vous avez collectivement accumulé 5 pénalités sur 16.} \\
\texttt{L'autre joueur a déplacé un objet sur sa grille.} \\
\\ 
\\ 
\texttt{Quelle est votre prochaine commande ?} \\
            }
        }
    }
    & \\ \\

    \theutterance \stepcounter{utterance}  
    & \multicolumn{4}{p{0.6\linewidth}}{
        \cellcolor[rgb]{0.9,0.9,0.9}{
            \makecell[{{p{\linewidth}}}]{
                \texttt{\tiny{[P1$\rangle$GM]}}
                \texttt{DÉPLACER: T, (6, 1)} \\
            }
        }
    }
    & & \\ \\

    \theutterance \stepcounter{utterance}  
    & & & \multicolumn{2}{p{0.3\linewidth}}{
        \cellcolor[rgb]{0.9,0.9,0.9}{
            \makecell[{{p{\linewidth}}}]{
                \texttt{\tiny{[GM$|$GM]}}
                \texttt{DÉPLACER: T, (6, 1)} \\
            }
        }
    }
    & & \\ \\

    \theutterance \stepcounter{utterance}  
    & & & \multicolumn{2}{p{0.3\linewidth}}{
        \cellcolor[rgb]{0.9,0.9,0.9}{
            \makecell[{{p{\linewidth}}}]{
                \texttt{\tiny{[GM$|$GM]}}
                \texttt{Pénalité : (6,1) n'est pas vide, mais contient '│'.} \\
\\ 
\texttt{Vous avez collectivement accumulé 6 pénalités sur 16.} \\
\texttt{Veuillez réessayer !} \\
            }
        }
    }
    & & \\ \\

    \theutterance \stepcounter{utterance}  
    & & \multicolumn{4}{p{0.6\linewidth}}{
        \cellcolor[rgb]{0.9,0.9,0.9}{
            \makecell[{{p{\linewidth}}}]{
                \texttt{\tiny{[P1$\langle$GM]}}
                \texttt{Pénalité : (6,1) n'est pas vide, mais contient '│'.} \\
\\ 
\texttt{Vous avez collectivement accumulé 6 pénalités sur 16.} \\
\texttt{Veuillez réessayer !} \\
            }
        }
    }
    & \\ \\

    \theutterance \stepcounter{utterance}  
    & \multicolumn{4}{p{0.6\linewidth}}{
        \cellcolor[rgb]{0.9,0.9,0.9}{
            \makecell[{{p{\linewidth}}}]{
                \texttt{\tiny{[P1$\rangle$GM]}}
                \texttt{DÉPLACER: T, (6, 2)} \\
            }
        }
    }
    & & \\ \\

    \theutterance \stepcounter{utterance}  
    & & & \multicolumn{2}{p{0.3\linewidth}}{
        \cellcolor[rgb]{0.9,0.9,0.9}{
            \makecell[{{p{\linewidth}}}]{
                \texttt{\tiny{[GM$|$GM]}}
                \texttt{DÉPLACER: T, (6, 2)} \\
            }
        }
    }
    & & \\ \\

    \theutterance \stepcounter{utterance}  
    & & & \multicolumn{2}{p{0.3\linewidth}}{
        \cellcolor[rgb]{0.9,0.9,0.9}{
            \makecell[{{p{\linewidth}}}]{
                \texttt{\tiny{[GM$|$GM]}}
                \texttt{Pénalité : (6,2) n'est pas vide, mais contient '┴'.} \\
\\ 
\texttt{Vous avez collectivement accumulé 7 pénalités sur 16.} \\
\texttt{Veuillez réessayer !} \\
            }
        }
    }
    & & \\ \\

    \theutterance \stepcounter{utterance}  
    & & \multicolumn{4}{p{0.6\linewidth}}{
        \cellcolor[rgb]{0.9,0.9,0.9}{
            \makecell[{{p{\linewidth}}}]{
                \texttt{\tiny{[P1$\langle$GM]}}
                \texttt{Pénalité : (6,2) n'est pas vide, mais contient '┴'.} \\
\\ 
\texttt{Vous avez collectivement accumulé 7 pénalités sur 16.} \\
\texttt{Veuillez réessayer !} \\
            }
        }
    }
    & \\ \\

    \theutterance \stepcounter{utterance}  
    & \multicolumn{4}{p{0.6\linewidth}}{
        \cellcolor[rgb]{0.9,0.9,0.9}{
            \makecell[{{p{\linewidth}}}]{
                \texttt{\tiny{[P1$\rangle$GM]}}
                \texttt{DÉPLACER: T, (7, 2)} \\
            }
        }
    }
    & & \\ \\

    \theutterance \stepcounter{utterance}  
    & & & \multicolumn{2}{p{0.3\linewidth}}{
        \cellcolor[rgb]{0.9,0.9,0.9}{
            \makecell[{{p{\linewidth}}}]{
                \texttt{\tiny{[GM$|$GM]}}
                \texttt{DÉPLACER: T, (7, 2)} \\
            }
        }
    }
    & & \\ \\

    \theutterance \stepcounter{utterance}  
    & & & \multicolumn{2}{p{0.3\linewidth}}{
        \cellcolor[rgb]{0.9,0.9,0.9}{
            \makecell[{{p{\linewidth}}}]{
                \texttt{\tiny{[GM$|$GM]}}
                \texttt{Pénalité : (7,2) n'est pas vide, mais contient '─'.} \\
\\ 
\texttt{Vous avez collectivement accumulé 8 pénalités sur 16.} \\
\texttt{Veuillez réessayer !} \\
            }
        }
    }
    & & \\ \\

    \theutterance \stepcounter{utterance}  
    & & \multicolumn{4}{p{0.6\linewidth}}{
        \cellcolor[rgb]{0.9,0.9,0.9}{
            \makecell[{{p{\linewidth}}}]{
                \texttt{\tiny{[P1$\langle$GM]}}
                \texttt{Pénalité : (7,2) n'est pas vide, mais contient '─'.} \\
\\ 
\texttt{Vous avez collectivement accumulé 8 pénalités sur 16.} \\
\texttt{Veuillez réessayer !} \\
            }
        }
    }
    & \\ \\

    \theutterance \stepcounter{utterance}  
    & \multicolumn{4}{p{0.6\linewidth}}{
        \cellcolor[rgb]{0.9,0.9,0.9}{
            \makecell[{{p{\linewidth}}}]{
                \texttt{\tiny{[P1$\rangle$GM]}}
                \texttt{DÉPLACER: T, (6, 3)} \\
            }
        }
    }
    & & \\ \\

    \theutterance \stepcounter{utterance}  
    & & & \multicolumn{2}{p{0.3\linewidth}}{
        \cellcolor[rgb]{0.9,0.9,0.9}{
            \makecell[{{p{\linewidth}}}]{
                \texttt{\tiny{[GM$|$GM]}}
                \texttt{DÉPLACER: T, (6, 3)} \\
            }
        }
    }
    & & \\ \\

    \theutterance \stepcounter{utterance}  
    & & & \multicolumn{2}{p{0.3\linewidth}}{
        \cellcolor[rgb]{0.9,0.9,0.9}{
            \makecell[{{p{\linewidth}}}]{
                \texttt{\tiny{[GM$|$GM]}}
                \texttt{T moved to (6, 3)} \\
            }
        }
    }
    & & \\ \\

    \theutterance \stepcounter{utterance}  
    & & \multicolumn{4}{p{0.6\linewidth}}{
        \cellcolor[rgb]{0.9,0.9,0.9}{
            \makecell[{{p{\linewidth}}}]{
                \texttt{\tiny{[GM$\rangle$P2]}}
                \texttt{Déplacement de 'T' vers (6,3) réussi. Votre grille mise à jour ressemble à ceci :} \\
\texttt{```} \\
\texttt{ 1234567} \\
\texttt{╔══╤══╤═╗} \\
\texttt{║◌◌│◌◌│◌║ 1} \\
\texttt{╟──┼──┴─╢ 2} \\
\texttt{║◌◌│◌STP║ 3} \\
\texttt{╟──┘◌◌MO║ 4} \\
\texttt{║D◌◌◌◌◌◌║ 5} \\
\texttt{║◌◌┌─┬──╢ 6} \\
\texttt{║◌◌│A│◌◌║ 7} \\
\texttt{╚══╧═╧══╝} \\
\texttt{```} \\
\texttt{Vous jouez actuellement le tour 4 sur un maximum de 28.} \\
\texttt{Vous avez collectivement accumulé 8 pénalités sur 16.} \\
\texttt{L'autre joueur a déplacé un objet sur sa grille.} \\
\\ 
\\ 
\texttt{Quelle est votre prochaine commande ?} \\
            }
        }
    }
    & \\ \\

    \theutterance \stepcounter{utterance}  
    & & & \multicolumn{4}{p{0.6\linewidth}}{
        \cellcolor[rgb]{0.9,0.9,0.9}{
            \makecell[{{p{\linewidth}}}]{
                \texttt{\tiny{[GM$\langle$P2]}}
                \texttt{DÉPLACER: P, (7, 1)} \\
            }
        }
    }
     \\ \\

    \theutterance \stepcounter{utterance}  
    & & & \multicolumn{2}{p{0.3\linewidth}}{
        \cellcolor[rgb]{0.9,0.9,0.9}{
            \makecell[{{p{\linewidth}}}]{
                \texttt{\tiny{[GM$|$GM]}}
                \texttt{DÉPLACER: P, (7, 1)} \\
            }
        }
    }
    & & \\ \\

    \theutterance \stepcounter{utterance}  
    & & & \multicolumn{2}{p{0.3\linewidth}}{
        \cellcolor[rgb]{0.9,0.9,0.9}{
            \makecell[{{p{\linewidth}}}]{
                \texttt{\tiny{[GM$|$GM]}}
                \texttt{P moved to (7, 1)} \\
            }
        }
    }
    & & \\ \\

    \theutterance \stepcounter{utterance}  
    & & \multicolumn{4}{p{0.6\linewidth}}{
        \cellcolor[rgb]{0.9,0.9,0.9}{
            \makecell[{{p{\linewidth}}}]{
                \texttt{\tiny{[P1$\langle$GM]}}
                \texttt{Déplacement de 'T' vers (6,3) réussi. Votre grille mise à jour ressemble à ceci :} \\
\texttt{```} \\
\texttt{ 1234567} \\
\texttt{╔══╤══╤═╗} \\
\texttt{║◌P│◌M│◌║ 1} \\
\texttt{╟──┼──┴─╢ 2} \\
\texttt{║◌◌│◌STO║ 3} \\
\texttt{╟──┘◌A◌◌║ 4} \\
\texttt{║D◌◌◌◌◌◌║ 5} \\
\texttt{║◌◌┌─┬──╢ 6} \\
\texttt{║◌◌│◌│◌◌║ 7} \\
\texttt{╚══╧═╧══╝} \\
\texttt{```} \\
\texttt{Vous jouez actuellement le tour 4 sur un maximum de 28.} \\
\texttt{Vous avez collectivement accumulé 8 pénalités sur 16.} \\
\texttt{L'autre joueur a déplacé un objet sur sa grille.} \\
\\ 
\\ 
\texttt{Quelle est votre prochaine commande ?} \\
            }
        }
    }
    & \\ \\

    \theutterance \stepcounter{utterance}  
    & \multicolumn{4}{p{0.6\linewidth}}{
        \cellcolor[rgb]{0.9,0.9,0.9}{
            \makecell[{{p{\linewidth}}}]{
                \texttt{\tiny{[P1$\rangle$GM]}}
                \texttt{DÉPLACER: P, (7, 1)} \\
            }
        }
    }
    & & \\ \\

    \theutterance \stepcounter{utterance}  
    & & & \multicolumn{2}{p{0.3\linewidth}}{
        \cellcolor[rgb]{0.9,0.9,0.9}{
            \makecell[{{p{\linewidth}}}]{
                \texttt{\tiny{[GM$|$GM]}}
                \texttt{DÉPLACER: P, (7, 1)} \\
            }
        }
    }
    & & \\ \\

    \theutterance \stepcounter{utterance}  
    & & & \multicolumn{2}{p{0.3\linewidth}}{
        \cellcolor[rgb]{0.9,0.9,0.9}{
            \makecell[{{p{\linewidth}}}]{
                \texttt{\tiny{[GM$|$GM]}}
                \texttt{P moved to (7, 1)} \\
            }
        }
    }
    & & \\ \\

    \theutterance \stepcounter{utterance}  
    & & \multicolumn{4}{p{0.6\linewidth}}{
        \cellcolor[rgb]{0.9,0.9,0.9}{
            \makecell[{{p{\linewidth}}}]{
                \texttt{\tiny{[GM$\rangle$P2]}}
                \texttt{Déplacement de 'P' vers (7,1) réussi. Votre grille mise à jour ressemble à ceci :} \\
\texttt{```} \\
\texttt{ 1234567} \\
\texttt{╔══╤══╤═╗} \\
\texttt{║◌◌│◌◌│P║ 1} \\
\texttt{╟──┼──┴─╢ 2} \\
\texttt{║◌◌│◌ST◌║ 3} \\
\texttt{╟──┘◌◌MO║ 4} \\
\texttt{║D◌◌◌◌◌◌║ 5} \\
\texttt{║◌◌┌─┬──╢ 6} \\
\texttt{║◌◌│A│◌◌║ 7} \\
\texttt{╚══╧═╧══╝} \\
\texttt{```} \\
\texttt{Vous jouez actuellement le tour 5 sur un maximum de 28.} \\
\texttt{Vous avez collectivement accumulé 8 pénalités sur 16.} \\
\texttt{L'autre joueur a déplacé un objet sur sa grille.} \\
\\ 
\\ 
\texttt{Quelle est votre prochaine commande ?} \\
            }
        }
    }
    & \\ \\

    \theutterance \stepcounter{utterance}  
    & & & \multicolumn{4}{p{0.6\linewidth}}{
        \cellcolor[rgb]{0.9,0.9,0.9}{
            \makecell[{{p{\linewidth}}}]{
                \texttt{\tiny{[GM$\langle$P2]}}
                \texttt{DÉPLACER: M, (6, 2)} \\
            }
        }
    }
     \\ \\

    \theutterance \stepcounter{utterance}  
    & & & \multicolumn{2}{p{0.3\linewidth}}{
        \cellcolor[rgb]{0.9,0.9,0.9}{
            \makecell[{{p{\linewidth}}}]{
                \texttt{\tiny{[GM$|$GM]}}
                \texttt{DÉPLACER: M, (6, 2)} \\
            }
        }
    }
    & & \\ \\

    \theutterance \stepcounter{utterance}  
    & & & \multicolumn{2}{p{0.3\linewidth}}{
        \cellcolor[rgb]{0.9,0.9,0.9}{
            \makecell[{{p{\linewidth}}}]{
                \texttt{\tiny{[GM$|$GM]}}
                \texttt{Pénalité : (6,2) n'est pas vide, mais contient '┴'.} \\
\\ 
\texttt{Vous avez collectivement accumulé 9 pénalités sur 16.} \\
\texttt{Veuillez réessayer !} \\
            }
        }
    }
    & & \\ \\

    \theutterance \stepcounter{utterance}  
    & & \multicolumn{4}{p{0.6\linewidth}}{
        \cellcolor[rgb]{0.9,0.9,0.9}{
            \makecell[{{p{\linewidth}}}]{
                \texttt{\tiny{[GM$\rangle$P2]}}
                \texttt{Pénalité : (6,2) n'est pas vide, mais contient '┴'.} \\
\\ 
\texttt{Vous avez collectivement accumulé 9 pénalités sur 16.} \\
\texttt{Veuillez réessayer !} \\
            }
        }
    }
    & \\ \\

    \theutterance \stepcounter{utterance}  
    & & & \multicolumn{4}{p{0.6\linewidth}}{
        \cellcolor[rgb]{0.9,0.9,0.9}{
            \makecell[{{p{\linewidth}}}]{
                \texttt{\tiny{[GM$\langle$P2]}}
                \texttt{DÉPLACER: M, (7, 2)} \\
            }
        }
    }
     \\ \\

    \theutterance \stepcounter{utterance}  
    & & & \multicolumn{2}{p{0.3\linewidth}}{
        \cellcolor[rgb]{0.9,0.9,0.9}{
            \makecell[{{p{\linewidth}}}]{
                \texttt{\tiny{[GM$|$GM]}}
                \texttt{DÉPLACER: M, (7, 2)} \\
            }
        }
    }
    & & \\ \\

    \theutterance \stepcounter{utterance}  
    & & & \multicolumn{2}{p{0.3\linewidth}}{
        \cellcolor[rgb]{0.9,0.9,0.9}{
            \makecell[{{p{\linewidth}}}]{
                \texttt{\tiny{[GM$|$GM]}}
                \texttt{Pénalité : (7,2) n'est pas vide, mais contient '─'.} \\
\\ 
\texttt{Vous avez collectivement accumulé 10 pénalités sur 16.} \\
\texttt{Veuillez réessayer !} \\
            }
        }
    }
    & & \\ \\

    \theutterance \stepcounter{utterance}  
    & & \multicolumn{4}{p{0.6\linewidth}}{
        \cellcolor[rgb]{0.9,0.9,0.9}{
            \makecell[{{p{\linewidth}}}]{
                \texttt{\tiny{[GM$\rangle$P2]}}
                \texttt{Pénalité : (7,2) n'est pas vide, mais contient '─'.} \\
\\ 
\texttt{Vous avez collectivement accumulé 10 pénalités sur 16.} \\
\texttt{Veuillez réessayer !} \\
            }
        }
    }
    & \\ \\

    \theutterance \stepcounter{utterance}  
    & & & \multicolumn{4}{p{0.6\linewidth}}{
        \cellcolor[rgb]{0.9,0.9,0.9}{
            \makecell[{{p{\linewidth}}}]{
                \texttt{\tiny{[GM$\langle$P2]}}
                \texttt{DÉPLACER: M, (6, 4)} \\
            }
        }
    }
     \\ \\

    \theutterance \stepcounter{utterance}  
    & & & \multicolumn{2}{p{0.3\linewidth}}{
        \cellcolor[rgb]{0.9,0.9,0.9}{
            \makecell[{{p{\linewidth}}}]{
                \texttt{\tiny{[GM$|$GM]}}
                \texttt{DÉPLACER: M, (6, 4)} \\
            }
        }
    }
    & & \\ \\

    \theutterance \stepcounter{utterance}  
    & & & \multicolumn{2}{p{0.3\linewidth}}{
        \cellcolor[rgb]{0.9,0.9,0.9}{
            \makecell[{{p{\linewidth}}}]{
                \texttt{\tiny{[GM$|$GM]}}
                \texttt{Pénalité : (6,4) n'est pas vide, mais contient 'M'.} \\
\\ 
\texttt{Vous avez collectivement accumulé 11 pénalités sur 16.} \\
\texttt{Veuillez réessayer !} \\
            }
        }
    }
    & & \\ \\

    \theutterance \stepcounter{utterance}  
    & & \multicolumn{4}{p{0.6\linewidth}}{
        \cellcolor[rgb]{0.9,0.9,0.9}{
            \makecell[{{p{\linewidth}}}]{
                \texttt{\tiny{[GM$\rangle$P2]}}
                \texttt{Pénalité : (6,4) n'est pas vide, mais contient 'M'.} \\
\\ 
\texttt{Vous avez collectivement accumulé 11 pénalités sur 16.} \\
\texttt{Veuillez réessayer !} \\
            }
        }
    }
    & \\ \\

    \theutterance \stepcounter{utterance}  
    & & & \multicolumn{4}{p{0.6\linewidth}}{
        \cellcolor[rgb]{0.9,0.9,0.9}{
            \makecell[{{p{\linewidth}}}]{
                \texttt{\tiny{[GM$\langle$P2]}}
                \texttt{DÉPLACER: M, (6, 1)} \\
            }
        }
    }
     \\ \\

    \theutterance \stepcounter{utterance}  
    & & & \multicolumn{2}{p{0.3\linewidth}}{
        \cellcolor[rgb]{0.9,0.9,0.9}{
            \makecell[{{p{\linewidth}}}]{
                \texttt{\tiny{[GM$|$GM]}}
                \texttt{DÉPLACER: M, (6, 1)} \\
            }
        }
    }
    & & \\ \\

    \theutterance \stepcounter{utterance}  
    & & & \multicolumn{2}{p{0.3\linewidth}}{
        \cellcolor[rgb]{0.9,0.9,0.9}{
            \makecell[{{p{\linewidth}}}]{
                \texttt{\tiny{[GM$|$GM]}}
                \texttt{Pénalité : (6,1) n'est pas vide, mais contient '│'.} \\
\\ 
\texttt{Vous avez collectivement accumulé 12 pénalités sur 16.} \\
\texttt{Veuillez réessayer !} \\
            }
        }
    }
    & & \\ \\

    \theutterance \stepcounter{utterance}  
    & & \multicolumn{4}{p{0.6\linewidth}}{
        \cellcolor[rgb]{0.9,0.9,0.9}{
            \makecell[{{p{\linewidth}}}]{
                \texttt{\tiny{[GM$\rangle$P2]}}
                \texttt{Pénalité : (6,1) n'est pas vide, mais contient '│'.} \\
\\ 
\texttt{Vous avez collectivement accumulé 12 pénalités sur 16.} \\
\texttt{Veuillez réessayer !} \\
            }
        }
    }
    & \\ \\

    \theutterance \stepcounter{utterance}  
    & & & \multicolumn{4}{p{0.6\linewidth}}{
        \cellcolor[rgb]{0.9,0.9,0.9}{
            \makecell[{{p{\linewidth}}}]{
                \texttt{\tiny{[GM$\langle$P2]}}
                \texttt{DÉPLACER: O, (7, 2)} \\
            }
        }
    }
     \\ \\

    \theutterance \stepcounter{utterance}  
    & & & \multicolumn{2}{p{0.3\linewidth}}{
        \cellcolor[rgb]{0.9,0.9,0.9}{
            \makecell[{{p{\linewidth}}}]{
                \texttt{\tiny{[GM$|$GM]}}
                \texttt{DÉPLACER: O, (7, 2)} \\
            }
        }
    }
    & & \\ \\

    \theutterance \stepcounter{utterance}  
    & & & \multicolumn{2}{p{0.3\linewidth}}{
        \cellcolor[rgb]{0.9,0.9,0.9}{
            \makecell[{{p{\linewidth}}}]{
                \texttt{\tiny{[GM$|$GM]}}
                \texttt{Pénalité : (7,2) n'est pas vide, mais contient '─'.} \\
\\ 
\texttt{Vous avez collectivement accumulé 13 pénalités sur 16.} \\
\texttt{Veuillez réessayer !} \\
            }
        }
    }
    & & \\ \\

    \theutterance \stepcounter{utterance}  
    & & \multicolumn{4}{p{0.6\linewidth}}{
        \cellcolor[rgb]{0.9,0.9,0.9}{
            \makecell[{{p{\linewidth}}}]{
                \texttt{\tiny{[GM$\rangle$P2]}}
                \texttt{Pénalité : (7,2) n'est pas vide, mais contient '─'.} \\
\\ 
\texttt{Vous avez collectivement accumulé 13 pénalités sur 16.} \\
\texttt{Veuillez réessayer !} \\
            }
        }
    }
    & \\ \\

    \theutterance \stepcounter{utterance}  
    & & & \multicolumn{4}{p{0.6\linewidth}}{
        \cellcolor[rgb]{0.9,0.9,0.9}{
            \makecell[{{p{\linewidth}}}]{
                \texttt{\tiny{[GM$\langle$P2]}}
                \texttt{DÉPLACER: O, (6, 4)} \\
            }
        }
    }
     \\ \\

    \theutterance \stepcounter{utterance}  
    & & & \multicolumn{2}{p{0.3\linewidth}}{
        \cellcolor[rgb]{0.9,0.9,0.9}{
            \makecell[{{p{\linewidth}}}]{
                \texttt{\tiny{[GM$|$GM]}}
                \texttt{DÉPLACER: O, (6, 4)} \\
            }
        }
    }
    & & \\ \\

    \theutterance \stepcounter{utterance}  
    & & & \multicolumn{2}{p{0.3\linewidth}}{
        \cellcolor[rgb]{0.9,0.9,0.9}{
            \makecell[{{p{\linewidth}}}]{
                \texttt{\tiny{[GM$|$GM]}}
                \texttt{Pénalité : (6,4) n'est pas vide, mais contient 'M'.} \\
\\ 
\texttt{Vous avez collectivement accumulé 14 pénalités sur 16.} \\
\texttt{Veuillez réessayer !} \\
            }
        }
    }
    & & \\ \\

    \theutterance \stepcounter{utterance}  
    & & \multicolumn{4}{p{0.6\linewidth}}{
        \cellcolor[rgb]{0.9,0.9,0.9}{
            \makecell[{{p{\linewidth}}}]{
                \texttt{\tiny{[GM$\rangle$P2]}}
                \texttt{Pénalité : (6,4) n'est pas vide, mais contient 'M'.} \\
\\ 
\texttt{Vous avez collectivement accumulé 14 pénalités sur 16.} \\
\texttt{Veuillez réessayer !} \\
            }
        }
    }
    & \\ \\

    \theutterance \stepcounter{utterance}  
    & & & \multicolumn{4}{p{0.6\linewidth}}{
        \cellcolor[rgb]{0.9,0.9,0.9}{
            \makecell[{{p{\linewidth}}}]{
                \texttt{\tiny{[GM$\langle$P2]}}
                \texttt{DÉPLACER: S, (7, 2)} \\
            }
        }
    }
     \\ \\

    \theutterance \stepcounter{utterance}  
    & & & \multicolumn{2}{p{0.3\linewidth}}{
        \cellcolor[rgb]{0.9,0.9,0.9}{
            \makecell[{{p{\linewidth}}}]{
                \texttt{\tiny{[GM$|$GM]}}
                \texttt{DÉPLACER: S, (7, 2)} \\
            }
        }
    }
    & & \\ \\

    \theutterance \stepcounter{utterance}  
    & & & \multicolumn{2}{p{0.3\linewidth}}{
        \cellcolor[rgb]{0.9,0.9,0.9}{
            \makecell[{{p{\linewidth}}}]{
                \texttt{\tiny{[GM$|$GM]}}
                \texttt{Pénalité : (7,2) n'est pas vide, mais contient '─'.} \\
\\ 
\texttt{Vous avez collectivement accumulé 15 pénalités sur 16.} \\
\texttt{Veuillez réessayer !} \\
            }
        }
    }
    & & \\ \\

    \theutterance \stepcounter{utterance}  
    & & \multicolumn{4}{p{0.6\linewidth}}{
        \cellcolor[rgb]{0.9,0.9,0.9}{
            \makecell[{{p{\linewidth}}}]{
                \texttt{\tiny{[GM$\rangle$P2]}}
                \texttt{Pénalité : (7,2) n'est pas vide, mais contient '─'.} \\
\\ 
\texttt{Vous avez collectivement accumulé 15 pénalités sur 16.} \\
\texttt{Veuillez réessayer !} \\
            }
        }
    }
    & \\ \\

    \theutterance \stepcounter{utterance}  
    & & & \multicolumn{4}{p{0.6\linewidth}}{
        \cellcolor[rgb]{0.9,0.9,0.9}{
            \makecell[{{p{\linewidth}}}]{
                \texttt{\tiny{[GM$\langle$P2]}}
                \texttt{DÉPLACER: S, (6, 4)} \\
            }
        }
    }
     \\ \\

    \theutterance \stepcounter{utterance}  
    & & & \multicolumn{2}{p{0.3\linewidth}}{
        \cellcolor[rgb]{0.9,0.9,0.9}{
            \makecell[{{p{\linewidth}}}]{
                \texttt{\tiny{[GM$|$GM]}}
                \texttt{DÉPLACER: S, (6, 4)} \\
            }
        }
    }
    & & \\ \\

    \theutterance \stepcounter{utterance}  
    & & & \multicolumn{2}{p{0.3\linewidth}}{
        \cellcolor[rgb]{0.9,0.9,0.9}{
            \makecell[{{p{\linewidth}}}]{
                \texttt{\tiny{[GM$|$GM]}}
                \texttt{Pénalité : (6,4) n'est pas vide, mais contient 'M'.} \\
\\ 
\texttt{Vous avez collectivement accumulé 16 pénalités sur 16.} \\
\texttt{Veuillez réessayer !} \\
            }
        }
    }
    & & \\ \\

    \theutterance \stepcounter{utterance}  
    & & \multicolumn{4}{p{0.6\linewidth}}{
        \cellcolor[rgb]{0.9,0.9,0.9}{
            \makecell[{{p{\linewidth}}}]{
                \texttt{\tiny{[GM$\rangle$P2]}}
                \texttt{Pénalité : (6,4) n'est pas vide, mais contient 'M'.} \\
\\ 
\texttt{Vous avez collectivement accumulé 16 pénalités sur 16.} \\
\texttt{Veuillez réessayer !} \\
            }
        }
    }
    & \\ \\

    \theutterance \stepcounter{utterance}  
    & & & \multicolumn{4}{p{0.6\linewidth}}{
        \cellcolor[rgb]{0.9,0.9,0.9}{
            \makecell[{{p{\linewidth}}}]{
                \texttt{\tiny{[GM$\langle$P2]}}
                \texttt{DIRE: La zone cible (6,1){-}(7,4) est bloquée par des murs (│, ─, ┴) sur ma grille. Impossible d'y placer les objets. On change de zone ? Le coin bas{-}droit (col 6{-}7, lig 5{-}7) a l'air libre : (6,5),(7,5),(6,6),(7,6),(6,7),(7,7). Ou le centre ? Dis{-}moi ce qui est libre chez toi.} \\
            }
        }
    }
     \\ \\

    \theutterance \stepcounter{utterance}  
    & & & \multicolumn{2}{p{0.3\linewidth}}{
        \cellcolor[rgb]{0.9,0.9,0.9}{
            \makecell[{{p{\linewidth}}}]{
                \texttt{\tiny{[GM$|$GM]}}
                \texttt{DIRE: La zone cible (6,1){-}(7,4) est bloquée par des murs (│, ─, ┴) sur ma grille. Impossible d'y placer les objets. On change de zone ? Le coin bas{-}droit (col 6{-}7, lig 5{-}7) a l'air libre : (6,5),(7,5),(6,6),(7,6),(6,7),(7,7). Ou le centre ? Dis{-}moi ce qui est libre chez toi.} \\
            }
        }
    }
    & & \\ \\

    \theutterance \stepcounter{utterance}  
    & & \multicolumn{4}{p{0.6\linewidth}}{
        \cellcolor[rgb]{0.9,0.9,0.9}{
            \makecell[{{p{\linewidth}}}]{
                \texttt{\tiny{[P1$\langle$GM]}}
                \texttt{Déplacement de 'P' vers (7,1) réussi. Votre grille mise à jour ressemble à ceci :} \\
\texttt{```} \\
\texttt{ 1234567} \\
\texttt{╔══╤══╤═╗} \\
\texttt{║◌◌│◌M│P║ 1} \\
\texttt{╟──┼──┴─╢ 2} \\
\texttt{║◌◌│◌STO║ 3} \\
\texttt{╟──┘◌A◌◌║ 4} \\
\texttt{║D◌◌◌◌◌◌║ 5} \\
\texttt{║◌◌┌─┬──╢ 6} \\
\texttt{║◌◌│◌│◌◌║ 7} \\
\texttt{╚══╧═╧══╝} \\
\texttt{```} \\
\texttt{Vous jouez actuellement le tour 5 sur un maximum de 28.} \\
\texttt{Vous avez collectivement accumulé 16 pénalités sur 16.} \\
\texttt{L'autre joueur a envoyé ce message :} \\
\\ 
\\ 
\texttt{"La zone cible (6,1){-}(7,4) est bloquée par des murs (│, ─, ┴) sur ma grille. Impossible d'y placer les objets. On change de zone ? Le coin bas{-}droit (col 6{-}7, lig 5{-}7) a l'air libre : (6,5),(7,5),(6,6),(7,6),(6,7),(7,7). Ou le centre ? Dis{-}moi ce qui est libre chez toi."} \\
\\ 
\\ 
\texttt{Quelle est votre prochaine commande ?} \\
            }
        }
    }
    & \\ \\

    \theutterance \stepcounter{utterance}  
    & \multicolumn{4}{p{0.6\linewidth}}{
        \cellcolor[rgb]{0.9,0.9,0.9}{
            \makecell[{{p{\linewidth}}}]{
                \texttt{\tiny{[P1$\rangle$GM]}}
                \texttt{DIRE: Le coin bas{-}droit (colonnes 6{-}7, lignes 5{-}7) est complètement vide chez moi. Allons{-}y. Nouvelle configuration : D(6,5), A(7,5), M(6,6), O(7,6), P(6,7), S(7,7), T(6,4) {[}T est déjà en (6,3), je le descendrai en (6,4){]}. Je libère d'abord (6,4) en déplaçant A de (3,4) vers (7,5).} \\
            }
        }
    }
    & & \\ \\

    \theutterance \stepcounter{utterance}  
    & & & \multicolumn{2}{p{0.3\linewidth}}{
        \cellcolor[rgb]{0.9,0.9,0.9}{
            \makecell[{{p{\linewidth}}}]{
                \texttt{\tiny{[GM$|$GM]}}
                \texttt{DIRE: Le coin bas{-}droit (colonnes 6{-}7, lignes 5{-}7) est complètement vide chez moi. Allons{-}y. Nouvelle configuration : D(6,5), A(7,5), M(6,6), O(7,6), P(6,7), S(7,7), T(6,4) {[}T est déjà en (6,3), je le descendrai en (6,4){]}. Je libère d'abord (6,4) en déplaçant A de (3,4) vers (7,5).} \\
            }
        }
    }
    & & \\ \\

    \theutterance \stepcounter{utterance}  
    & & \multicolumn{4}{p{0.6\linewidth}}{
        \cellcolor[rgb]{0.9,0.9,0.9}{
            \makecell[{{p{\linewidth}}}]{
                \texttt{\tiny{[GM$\rangle$P2]}}
                \texttt{Votre message a été transmis à l'autre joueur.} \\
\\ 
\\ 
\texttt{Vous jouez actuellement le tour 6 sur un maximum de 28.} \\
\texttt{Vous avez collectivement accumulé 16 pénalités sur 16.} \\
\texttt{L'autre joueur a envoyé ce message :} \\
\\ 
\\ 
\texttt{"Le coin bas{-}droit (colonnes 6{-}7, lignes 5{-}7) est complètement vide chez moi. Allons{-}y. Nouvelle configuration : D(6,5), A(7,5), M(6,6), O(7,6), P(6,7), S(7,7), T(6,4) {[}T est déjà en (6,3), je le descendrai en (6,4){]}. Je libère d'abord (6,4) en déplaçant A de (3,4) vers (7,5)."} \\
\\ 
\\ 
\texttt{Quelle est votre prochaine commande ?} \\
            }
        }
    }
    & \\ \\

    \theutterance \stepcounter{utterance}  
    & & & \multicolumn{4}{p{0.6\linewidth}}{
        \cellcolor[rgb]{0.9,0.9,0.9}{
            \makecell[{{p{\linewidth}}}]{
                \texttt{\tiny{[GM$\langle$P2]}}
                \texttt{DÉPLACER: T, (6, 4)} \\
            }
        }
    }
     \\ \\

    \theutterance \stepcounter{utterance}  
    & & & \multicolumn{2}{p{0.3\linewidth}}{
        \cellcolor[rgb]{0.9,0.9,0.9}{
            \makecell[{{p{\linewidth}}}]{
                \texttt{\tiny{[GM$|$GM]}}
                \texttt{DÉPLACER: T, (6, 4)} \\
            }
        }
    }
    & & \\ \\

    \theutterance \stepcounter{utterance}  
    & & & \multicolumn{2}{p{0.3\linewidth}}{
        \cellcolor[rgb]{0.9,0.9,0.9}{
            \makecell[{{p{\linewidth}}}]{
                \texttt{\tiny{[GM$|$GM]}}
                \texttt{Pénalité : (6,4) n'est pas vide, mais contient 'M'.} \\
\\ 
\texttt{Vous avez collectivement accumulé 17 pénalités sur 16.} \\
\texttt{Veuillez réessayer !} \\
            }
        }
    }
    & & \\ \\

    \theutterance \stepcounter{utterance}  
    & & \multicolumn{4}{p{0.6\linewidth}}{
        \cellcolor[rgb]{0.9,0.9,0.9}{
            \makecell[{{p{\linewidth}}}]{
                \texttt{\tiny{[GM$\rangle$P2]}}
                \texttt{Pénalité : (6,4) n'est pas vide, mais contient 'M'.} \\
\\ 
\texttt{Vous avez collectivement accumulé 17 pénalités sur 16.} \\
\texttt{Veuillez réessayer !} \\
            }
        }
    }
    & \\ \\

    \theutterance \stepcounter{utterance}  
    & & & \multicolumn{4}{p{0.6\linewidth}}{
        \cellcolor[rgb]{0.9,0.9,0.9}{
            \makecell[{{p{\linewidth}}}]{
                \texttt{\tiny{[GM$\langle$P2]}}
                \texttt{DÉPLACER: M, (6, 5)} \\
            }
        }
    }
     \\ \\

    \theutterance \stepcounter{utterance}  
    & & & \multicolumn{2}{p{0.3\linewidth}}{
        \cellcolor[rgb]{0.9,0.9,0.9}{
            \makecell[{{p{\linewidth}}}]{
                \texttt{\tiny{[GM$|$GM]}}
                \texttt{DÉPLACER: M, (6, 5)} \\
            }
        }
    }
    & & \\ \\

    \theutterance \stepcounter{utterance}  
    & & & \multicolumn{2}{p{0.3\linewidth}}{
        \cellcolor[rgb]{0.9,0.9,0.9}{
            \makecell[{{p{\linewidth}}}]{
                \texttt{\tiny{[GM$|$GM]}}
                \texttt{M moved to (6, 5)} \\
            }
        }
    }
    & & \\ \\

    \theutterance \stepcounter{utterance}  
    & & \multicolumn{4}{p{0.6\linewidth}}{
        \cellcolor[rgb]{0.9,0.9,0.9}{
            \makecell[{{p{\linewidth}}}]{
                \texttt{\tiny{[P1$\langle$GM]}}
                \texttt{Votre message a été transmis à l'autre joueur.} \\
\\ 
\\ 
\texttt{Vous jouez actuellement le tour 6 sur un maximum de 28.} \\
\texttt{Vous avez collectivement accumulé 17 pénalités sur 16.} \\
\texttt{L'autre joueur a déplacé un objet sur sa grille.} \\
\\ 
\\ 
\texttt{Quelle est votre prochaine commande ?} \\
            }
        }
    }
    & \\ \\

    \theutterance \stepcounter{utterance}  
    & \multicolumn{4}{p{0.6\linewidth}}{
        \cellcolor[rgb]{0.9,0.9,0.9}{
            \makecell[{{p{\linewidth}}}]{
                \texttt{\tiny{[P1$\rangle$GM]}}
                \texttt{DÉPLACER: A, (7, 5)} \\
            }
        }
    }
    & & \\ \\

    \theutterance \stepcounter{utterance}  
    & & & \multicolumn{2}{p{0.3\linewidth}}{
        \cellcolor[rgb]{0.9,0.9,0.9}{
            \makecell[{{p{\linewidth}}}]{
                \texttt{\tiny{[GM$|$GM]}}
                \texttt{DÉPLACER: A, (7, 5)} \\
            }
        }
    }
    & & \\ \\

    \theutterance \stepcounter{utterance}  
    & & & \multicolumn{2}{p{0.3\linewidth}}{
        \cellcolor[rgb]{0.9,0.9,0.9}{
            \makecell[{{p{\linewidth}}}]{
                \texttt{\tiny{[GM$|$GM]}}
                \texttt{A moved to (7, 5)} \\
            }
        }
    }
    & & \\ \\

    \theutterance \stepcounter{utterance}  
    & & \multicolumn{4}{p{0.6\linewidth}}{
        \cellcolor[rgb]{0.9,0.9,0.9}{
            \makecell[{{p{\linewidth}}}]{
                \texttt{\tiny{[GM$\rangle$P2]}}
                \texttt{Déplacement de 'M' vers (6,5) réussi. Votre grille mise à jour ressemble à ceci :} \\
\texttt{```} \\
\texttt{ 1234567} \\
\texttt{╔══╤══╤═╗} \\
\texttt{║◌◌│◌◌│P║ 1} \\
\texttt{╟──┼──┴─╢ 2} \\
\texttt{║◌◌│◌ST◌║ 3} \\
\texttt{╟──┘◌◌◌O║ 4} \\
\texttt{║D◌◌◌◌M◌║ 5} \\
\texttt{║◌◌┌─┬──╢ 6} \\
\texttt{║◌◌│A│◌◌║ 7} \\
\texttt{╚══╧═╧══╝} \\
\texttt{```} \\
\texttt{Vous jouez actuellement le tour 7 sur un maximum de 28.} \\
\texttt{Vous avez collectivement accumulé 17 pénalités sur 16.} \\
\texttt{L'autre joueur a déplacé un objet sur sa grille.} \\
\\ 
\\ 
\texttt{Quelle est votre prochaine commande ?} \\
            }
        }
    }
    & \\ \\

    \theutterance \stepcounter{utterance}  
    & & & \multicolumn{4}{p{0.6\linewidth}}{
        \cellcolor[rgb]{0.9,0.9,0.9}{
            \makecell[{{p{\linewidth}}}]{
                \texttt{\tiny{[GM$\langle$P2]}}
                \texttt{DÉPLACER: T, (6, 4)} \\
            }
        }
    }
     \\ \\

    \theutterance \stepcounter{utterance}  
    & & & \multicolumn{2}{p{0.3\linewidth}}{
        \cellcolor[rgb]{0.9,0.9,0.9}{
            \makecell[{{p{\linewidth}}}]{
                \texttt{\tiny{[GM$|$GM]}}
                \texttt{DÉPLACER: T, (6, 4)} \\
            }
        }
    }
    & & \\ \\

    \theutterance \stepcounter{utterance}  
    & & & \multicolumn{2}{p{0.3\linewidth}}{
        \cellcolor[rgb]{0.9,0.9,0.9}{
            \makecell[{{p{\linewidth}}}]{
                \texttt{\tiny{[GM$|$GM]}}
                \texttt{T moved to (6, 4)} \\
            }
        }
    }
    & & \\ \\

    \theutterance \stepcounter{utterance}  
    & & \multicolumn{4}{p{0.6\linewidth}}{
        \cellcolor[rgb]{0.9,0.9,0.9}{
            \makecell[{{p{\linewidth}}}]{
                \texttt{\tiny{[P1$\langle$GM]}}
                \texttt{Déplacement de 'A' vers (7,5) réussi. Votre grille mise à jour ressemble à ceci :} \\
\texttt{```} \\
\texttt{ 1234567} \\
\texttt{╔══╤══╤═╗} \\
\texttt{║◌◌│◌M│P║ 1} \\
\texttt{╟──┼──┴─╢ 2} \\
\texttt{║◌◌│◌STO║ 3} \\
\texttt{╟──┘◌◌◌◌║ 4} \\
\texttt{║D◌◌◌◌◌A║ 5} \\
\texttt{║◌◌┌─┬──╢ 6} \\
\texttt{║◌◌│◌│◌◌║ 7} \\
\texttt{╚══╧═╧══╝} \\
\texttt{```} \\
\texttt{Vous jouez actuellement le tour 7 sur un maximum de 28.} \\
\texttt{Vous avez collectivement accumulé 17 pénalités sur 16.} \\
\texttt{L'autre joueur a déplacé un objet sur sa grille.} \\
\\ 
\\ 
\texttt{Quelle est votre prochaine commande ?} \\
            }
        }
    }
    & \\ \\

    \theutterance \stepcounter{utterance}  
    & \multicolumn{4}{p{0.6\linewidth}}{
        \cellcolor[rgb]{0.9,0.9,0.9}{
            \makecell[{{p{\linewidth}}}]{
                \texttt{\tiny{[P1$\rangle$GM]}}
                \texttt{DÉPLACER: D, (6, 5)} \\
            }
        }
    }
    & & \\ \\

    \theutterance \stepcounter{utterance}  
    & & & \multicolumn{2}{p{0.3\linewidth}}{
        \cellcolor[rgb]{0.9,0.9,0.9}{
            \makecell[{{p{\linewidth}}}]{
                \texttt{\tiny{[GM$|$GM]}}
                \texttt{DÉPLACER: D, (6, 5)} \\
            }
        }
    }
    & & \\ \\

    \theutterance \stepcounter{utterance}  
    & & & \multicolumn{2}{p{0.3\linewidth}}{
        \cellcolor[rgb]{0.9,0.9,0.9}{
            \makecell[{{p{\linewidth}}}]{
                \texttt{\tiny{[GM$|$GM]}}
                \texttt{D moved to (6, 5)} \\
            }
        }
    }
    & & \\ \\

    \theutterance \stepcounter{utterance}  
    & & \multicolumn{4}{p{0.6\linewidth}}{
        \cellcolor[rgb]{0.9,0.9,0.9}{
            \makecell[{{p{\linewidth}}}]{
                \texttt{\tiny{[GM$\rangle$P2]}}
                \texttt{Déplacement de 'T' vers (6,4) réussi. Votre grille mise à jour ressemble à ceci :} \\
\texttt{```} \\
\texttt{ 1234567} \\
\texttt{╔══╤══╤═╗} \\
\texttt{║◌◌│◌◌│P║ 1} \\
\texttt{╟──┼──┴─╢ 2} \\
\texttt{║◌◌│◌S◌◌║ 3} \\
\texttt{╟──┘◌◌TO║ 4} \\
\texttt{║D◌◌◌◌M◌║ 5} \\
\texttt{║◌◌┌─┬──╢ 6} \\
\texttt{║◌◌│A│◌◌║ 7} \\
\texttt{╚══╧═╧══╝} \\
\texttt{```} \\
\texttt{Vous jouez actuellement le tour 8 sur un maximum de 28.} \\
\texttt{Vous avez collectivement accumulé 17 pénalités sur 16.} \\
\texttt{L'autre joueur a déplacé un objet sur sa grille.} \\
\\ 
\\ 
\texttt{Quelle est votre prochaine commande ?} \\
            }
        }
    }
    & \\ \\

    \theutterance \stepcounter{utterance}  
    & & & \multicolumn{4}{p{0.6\linewidth}}{
        \cellcolor[rgb]{0.9,0.9,0.9}{
            \makecell[{{p{\linewidth}}}]{
                \texttt{\tiny{[GM$\langle$P2]}}
                \texttt{DÉPLACER: O, (7, 5)} \\
            }
        }
    }
     \\ \\

    \theutterance \stepcounter{utterance}  
    & & & \multicolumn{2}{p{0.3\linewidth}}{
        \cellcolor[rgb]{0.9,0.9,0.9}{
            \makecell[{{p{\linewidth}}}]{
                \texttt{\tiny{[GM$|$GM]}}
                \texttt{DÉPLACER: O, (7, 5)} \\
            }
        }
    }
    & & \\ \\

    \theutterance \stepcounter{utterance}  
    & & & \multicolumn{2}{p{0.3\linewidth}}{
        \cellcolor[rgb]{0.9,0.9,0.9}{
            \makecell[{{p{\linewidth}}}]{
                \texttt{\tiny{[GM$|$GM]}}
                \texttt{O moved to (7, 5)} \\
            }
        }
    }
    & & \\ \\

    \theutterance \stepcounter{utterance}  
    & & \multicolumn{4}{p{0.6\linewidth}}{
        \cellcolor[rgb]{0.9,0.9,0.9}{
            \makecell[{{p{\linewidth}}}]{
                \texttt{\tiny{[P1$\langle$GM]}}
                \texttt{Déplacement de 'D' vers (6,5) réussi. Votre grille mise à jour ressemble à ceci :} \\
\texttt{```} \\
\texttt{ 1234567} \\
\texttt{╔══╤══╤═╗} \\
\texttt{║◌◌│◌M│P║ 1} \\
\texttt{╟──┼──┴─╢ 2} \\
\texttt{║◌◌│◌STO║ 3} \\
\texttt{╟──┘◌◌◌◌║ 4} \\
\texttt{║◌◌◌◌◌DA║ 5} \\
\texttt{║◌◌┌─┬──╢ 6} \\
\texttt{║◌◌│◌│◌◌║ 7} \\
\texttt{╚══╧═╧══╝} \\
\texttt{```} \\
\texttt{Vous jouez actuellement le tour 8 sur un maximum de 28.} \\
\texttt{Vous avez collectivement accumulé 17 pénalités sur 16.} \\
\texttt{L'autre joueur a déplacé un objet sur sa grille.} \\
\\ 
\\ 
\texttt{Quelle est votre prochaine commande ?} \\
            }
        }
    }
    & \\ \\

    \theutterance \stepcounter{utterance}  
    & \multicolumn{4}{p{0.6\linewidth}}{
        \cellcolor[rgb]{0.9,0.9,0.9}{
            \makecell[{{p{\linewidth}}}]{
                \texttt{\tiny{[P1$\rangle$GM]}}
                \texttt{DÉPLACER: M, (6, 6)} \\
            }
        }
    }
    & & \\ \\

    \theutterance \stepcounter{utterance}  
    & & & \multicolumn{2}{p{0.3\linewidth}}{
        \cellcolor[rgb]{0.9,0.9,0.9}{
            \makecell[{{p{\linewidth}}}]{
                \texttt{\tiny{[GM$|$GM]}}
                \texttt{DÉPLACER: M, (6, 6)} \\
            }
        }
    }
    & & \\ \\

    \theutterance \stepcounter{utterance}  
    & & & \multicolumn{2}{p{0.3\linewidth}}{
        \cellcolor[rgb]{0.9,0.9,0.9}{
            \makecell[{{p{\linewidth}}}]{
                \texttt{\tiny{[GM$|$GM]}}
                \texttt{Pénalité : (6,6) n'est pas vide, mais contient '─'.} \\
\\ 
\texttt{Vous avez collectivement accumulé 18 pénalités sur 16.} \\
\texttt{Veuillez réessayer !} \\
            }
        }
    }
    & & \\ \\

    \theutterance \stepcounter{utterance}  
    & & \multicolumn{4}{p{0.6\linewidth}}{
        \cellcolor[rgb]{0.9,0.9,0.9}{
            \makecell[{{p{\linewidth}}}]{
                \texttt{\tiny{[P1$\langle$GM]}}
                \texttt{Pénalité : (6,6) n'est pas vide, mais contient '─'.} \\
\\ 
\texttt{Vous avez collectivement accumulé 18 pénalités sur 16.} \\
\texttt{Veuillez réessayer !} \\
            }
        }
    }
    & \\ \\

    \theutterance \stepcounter{utterance}  
    & \multicolumn{4}{p{0.6\linewidth}}{
        \cellcolor[rgb]{0.9,0.9,0.9}{
            \makecell[{{p{\linewidth}}}]{
                \texttt{\tiny{[P1$\rangle$GM]}}
                \texttt{DÉPLACER: M, (7, 6)} \\
            }
        }
    }
    & & \\ \\

    \theutterance \stepcounter{utterance}  
    & & & \multicolumn{2}{p{0.3\linewidth}}{
        \cellcolor[rgb]{0.9,0.9,0.9}{
            \makecell[{{p{\linewidth}}}]{
                \texttt{\tiny{[GM$|$GM]}}
                \texttt{DÉPLACER: M, (7, 6)} \\
            }
        }
    }
    & & \\ \\

    \theutterance \stepcounter{utterance}  
    & & & \multicolumn{2}{p{0.3\linewidth}}{
        \cellcolor[rgb]{0.9,0.9,0.9}{
            \makecell[{{p{\linewidth}}}]{
                \texttt{\tiny{[GM$|$GM]}}
                \texttt{Pénalité : (7,6) n'est pas vide, mais contient '─'.} \\
\\ 
\texttt{Vous avez collectivement accumulé 19 pénalités sur 16.} \\
\texttt{Veuillez réessayer !} \\
            }
        }
    }
    & & \\ \\

    \theutterance \stepcounter{utterance}  
    & & \multicolumn{4}{p{0.6\linewidth}}{
        \cellcolor[rgb]{0.9,0.9,0.9}{
            \makecell[{{p{\linewidth}}}]{
                \texttt{\tiny{[P1$\langle$GM]}}
                \texttt{Pénalité : (7,6) n'est pas vide, mais contient '─'.} \\
\\ 
\texttt{Vous avez collectivement accumulé 19 pénalités sur 16.} \\
\texttt{Veuillez réessayer !} \\
            }
        }
    }
    & \\ \\

    \theutterance \stepcounter{utterance}  
    & \multicolumn{4}{p{0.6\linewidth}}{
        \cellcolor[rgb]{0.9,0.9,0.9}{
            \makecell[{{p{\linewidth}}}]{
                \texttt{\tiny{[P1$\rangle$GM]}}
                \texttt{DÉPLACER: M, (6, 7)} \\
            }
        }
    }
    & & \\ \\

    \theutterance \stepcounter{utterance}  
    & & & \multicolumn{2}{p{0.3\linewidth}}{
        \cellcolor[rgb]{0.9,0.9,0.9}{
            \makecell[{{p{\linewidth}}}]{
                \texttt{\tiny{[GM$|$GM]}}
                \texttt{DÉPLACER: M, (6, 7)} \\
            }
        }
    }
    & & \\ \\

    \theutterance \stepcounter{utterance}  
    & & & \multicolumn{2}{p{0.3\linewidth}}{
        \cellcolor[rgb]{0.9,0.9,0.9}{
            \makecell[{{p{\linewidth}}}]{
                \texttt{\tiny{[GM$|$GM]}}
                \texttt{M moved to (6, 7)} \\
            }
        }
    }
    & & \\ \\

    \theutterance \stepcounter{utterance}  
    & & \multicolumn{4}{p{0.6\linewidth}}{
        \cellcolor[rgb]{0.9,0.9,0.9}{
            \makecell[{{p{\linewidth}}}]{
                \texttt{\tiny{[GM$\rangle$P2]}}
                \texttt{Déplacement de 'O' vers (7,5) réussi. Votre grille mise à jour ressemble à ceci :} \\
\texttt{```} \\
\texttt{ 1234567} \\
\texttt{╔══╤══╤═╗} \\
\texttt{║◌◌│◌◌│P║ 1} \\
\texttt{╟──┼──┴─╢ 2} \\
\texttt{║◌◌│◌S◌◌║ 3} \\
\texttt{╟──┘◌◌T◌║ 4} \\
\texttt{║D◌◌◌◌MO║ 5} \\
\texttt{║◌◌┌─┬──╢ 6} \\
\texttt{║◌◌│A│◌◌║ 7} \\
\texttt{╚══╧═╧══╝} \\
\texttt{```} \\
\texttt{Vous jouez actuellement le tour 9 sur un maximum de 28.} \\
\texttt{Vous avez collectivement accumulé 19 pénalités sur 16.} \\
\texttt{L'autre joueur a déplacé un objet sur sa grille.} \\
\\ 
\\ 
\texttt{Quelle est votre prochaine commande ?} \\
            }
        }
    }
    & \\ \\

    \theutterance \stepcounter{utterance}  
    & & & \multicolumn{4}{p{0.6\linewidth}}{
        \cellcolor[rgb]{0.9,0.9,0.9}{
            \makecell[{{p{\linewidth}}}]{
                \texttt{\tiny{[GM$\langle$P2]}}
                \texttt{DÉPLACER: S, (6, 6)} \\
            }
        }
    }
     \\ \\

    \theutterance \stepcounter{utterance}  
    & & & \multicolumn{2}{p{0.3\linewidth}}{
        \cellcolor[rgb]{0.9,0.9,0.9}{
            \makecell[{{p{\linewidth}}}]{
                \texttt{\tiny{[GM$|$GM]}}
                \texttt{DÉPLACER: S, (6, 6)} \\
            }
        }
    }
    & & \\ \\

    \theutterance \stepcounter{utterance}  
    & & & \multicolumn{2}{p{0.3\linewidth}}{
        \cellcolor[rgb]{0.9,0.9,0.9}{
            \makecell[{{p{\linewidth}}}]{
                \texttt{\tiny{[GM$|$GM]}}
                \texttt{Pénalité : (6,6) n'est pas vide, mais contient '─'.} \\
\\ 
\texttt{Vous avez collectivement accumulé 20 pénalités sur 16.} \\
\texttt{Veuillez réessayer !} \\
            }
        }
    }
    & & \\ \\

    \theutterance \stepcounter{utterance}  
    & & \multicolumn{4}{p{0.6\linewidth}}{
        \cellcolor[rgb]{0.9,0.9,0.9}{
            \makecell[{{p{\linewidth}}}]{
                \texttt{\tiny{[GM$\rangle$P2]}}
                \texttt{Pénalité : (6,6) n'est pas vide, mais contient '─'.} \\
\\ 
\texttt{Vous avez collectivement accumulé 20 pénalités sur 16.} \\
\texttt{Veuillez réessayer !} \\
            }
        }
    }
    & \\ \\

    \theutterance \stepcounter{utterance}  
    & & & \multicolumn{4}{p{0.6\linewidth}}{
        \cellcolor[rgb]{0.9,0.9,0.9}{
            \makecell[{{p{\linewidth}}}]{
                \texttt{\tiny{[GM$\langle$P2]}}
                \texttt{DÉPLACER: S, (7, 6)} \\
            }
        }
    }
     \\ \\

    \theutterance \stepcounter{utterance}  
    & & & \multicolumn{2}{p{0.3\linewidth}}{
        \cellcolor[rgb]{0.9,0.9,0.9}{
            \makecell[{{p{\linewidth}}}]{
                \texttt{\tiny{[GM$|$GM]}}
                \texttt{DÉPLACER: S, (7, 6)} \\
            }
        }
    }
    & & \\ \\

    \theutterance \stepcounter{utterance}  
    & & & \multicolumn{2}{p{0.3\linewidth}}{
        \cellcolor[rgb]{0.9,0.9,0.9}{
            \makecell[{{p{\linewidth}}}]{
                \texttt{\tiny{[GM$|$GM]}}
                \texttt{Pénalité : (7,6) n'est pas vide, mais contient '─'.} \\
\\ 
\texttt{Vous avez collectivement accumulé 21 pénalités sur 16.} \\
\texttt{Veuillez réessayer !} \\
            }
        }
    }
    & & \\ \\

    \theutterance \stepcounter{utterance}  
    & & \multicolumn{4}{p{0.6\linewidth}}{
        \cellcolor[rgb]{0.9,0.9,0.9}{
            \makecell[{{p{\linewidth}}}]{
                \texttt{\tiny{[GM$\rangle$P2]}}
                \texttt{Pénalité : (7,6) n'est pas vide, mais contient '─'.} \\
\\ 
\texttt{Vous avez collectivement accumulé 21 pénalités sur 16.} \\
\texttt{Veuillez réessayer !} \\
            }
        }
    }
    & \\ \\

    \theutterance \stepcounter{utterance}  
    & & & \multicolumn{4}{p{0.6\linewidth}}{
        \cellcolor[rgb]{0.9,0.9,0.9}{
            \makecell[{{p{\linewidth}}}]{
                \texttt{\tiny{[GM$\langle$P2]}}
                \texttt{DÉPLACER: S, (6, 7)} \\
            }
        }
    }
     \\ \\

    \theutterance \stepcounter{utterance}  
    & & & \multicolumn{2}{p{0.3\linewidth}}{
        \cellcolor[rgb]{0.9,0.9,0.9}{
            \makecell[{{p{\linewidth}}}]{
                \texttt{\tiny{[GM$|$GM]}}
                \texttt{DÉPLACER: S, (6, 7)} \\
            }
        }
    }
    & & \\ \\

    \theutterance \stepcounter{utterance}  
    & & & \multicolumn{2}{p{0.3\linewidth}}{
        \cellcolor[rgb]{0.9,0.9,0.9}{
            \makecell[{{p{\linewidth}}}]{
                \texttt{\tiny{[GM$|$GM]}}
                \texttt{S moved to (6, 7)} \\
            }
        }
    }
    & & \\ \\

    \theutterance \stepcounter{utterance}  
    & & \multicolumn{4}{p{0.6\linewidth}}{
        \cellcolor[rgb]{0.9,0.9,0.9}{
            \makecell[{{p{\linewidth}}}]{
                \texttt{\tiny{[P1$\langle$GM]}}
                \texttt{Déplacement de 'M' vers (6,7) réussi. Votre grille mise à jour ressemble à ceci :} \\
\texttt{```} \\
\texttt{ 1234567} \\
\texttt{╔══╤══╤═╗} \\
\texttt{║◌◌│◌◌│P║ 1} \\
\texttt{╟──┼──┴─╢ 2} \\
\texttt{║◌◌│◌STO║ 3} \\
\texttt{╟──┘◌◌◌◌║ 4} \\
\texttt{║◌◌◌◌◌DA║ 5} \\
\texttt{║◌◌┌─┬──╢ 6} \\
\texttt{║◌◌│◌│M◌║ 7} \\
\texttt{╚══╧═╧══╝} \\
\texttt{```} \\
\texttt{Vous jouez actuellement le tour 9 sur un maximum de 28.} \\
\texttt{Vous avez collectivement accumulé 21 pénalités sur 16.} \\
\texttt{L'autre joueur a déplacé un objet sur sa grille.} \\
\\ 
\\ 
\texttt{Quelle est votre prochaine commande ?} \\
            }
        }
    }
    & \\ \\

    \theutterance \stepcounter{utterance}  
    & \multicolumn{4}{p{0.6\linewidth}}{
        \cellcolor[rgb]{0.9,0.9,0.9}{
            \makecell[{{p{\linewidth}}}]{
                \texttt{\tiny{[P1$\rangle$GM]}}
                \texttt{DÉPLACER: O, (7, 7)} \\
            }
        }
    }
    & & \\ \\

    \theutterance \stepcounter{utterance}  
    & & & \multicolumn{2}{p{0.3\linewidth}}{
        \cellcolor[rgb]{0.9,0.9,0.9}{
            \makecell[{{p{\linewidth}}}]{
                \texttt{\tiny{[GM$|$GM]}}
                \texttt{DÉPLACER: O, (7, 7)} \\
            }
        }
    }
    & & \\ \\

    \theutterance \stepcounter{utterance}  
    & & & \multicolumn{2}{p{0.3\linewidth}}{
        \cellcolor[rgb]{0.9,0.9,0.9}{
            \makecell[{{p{\linewidth}}}]{
                \texttt{\tiny{[GM$|$GM]}}
                \texttt{O moved to (7, 7)} \\
            }
        }
    }
    & & \\ \\

    \theutterance \stepcounter{utterance}  
    & & \multicolumn{4}{p{0.6\linewidth}}{
        \cellcolor[rgb]{0.9,0.9,0.9}{
            \makecell[{{p{\linewidth}}}]{
                \texttt{\tiny{[GM$\rangle$P2]}}
                \texttt{Déplacement de 'S' vers (6,7) réussi. Votre grille mise à jour ressemble à ceci :} \\
\texttt{```} \\
\texttt{ 1234567} \\
\texttt{╔══╤══╤═╗} \\
\texttt{║◌◌│◌◌│P║ 1} \\
\texttt{╟──┼──┴─╢ 2} \\
\texttt{║◌◌│◌◌◌◌║ 3} \\
\texttt{╟──┘◌◌T◌║ 4} \\
\texttt{║D◌◌◌◌MO║ 5} \\
\texttt{║◌◌┌─┬──╢ 6} \\
\texttt{║◌◌│A│S◌║ 7} \\
\texttt{╚══╧═╧══╝} \\
\texttt{```} \\
\texttt{Vous jouez actuellement le tour 10 sur un maximum de 28.} \\
\texttt{Vous avez collectivement accumulé 21 pénalités sur 16.} \\
\texttt{L'autre joueur a déplacé un objet sur sa grille.} \\
\\ 
\\ 
\texttt{Quelle est votre prochaine commande ?} \\
            }
        }
    }
    & \\ \\

    \theutterance \stepcounter{utterance}  
    & & & \multicolumn{4}{p{0.6\linewidth}}{
        \cellcolor[rgb]{0.9,0.9,0.9}{
            \makecell[{{p{\linewidth}}}]{
                \texttt{\tiny{[GM$\langle$P2]}}
                \texttt{DÉPLACER: P, (7, 7)} \\
            }
        }
    }
     \\ \\

    \theutterance \stepcounter{utterance}  
    & & & \multicolumn{2}{p{0.3\linewidth}}{
        \cellcolor[rgb]{0.9,0.9,0.9}{
            \makecell[{{p{\linewidth}}}]{
                \texttt{\tiny{[GM$|$GM]}}
                \texttt{DÉPLACER: P, (7, 7)} \\
            }
        }
    }
    & & \\ \\

    \theutterance \stepcounter{utterance}  
    & & & \multicolumn{2}{p{0.3\linewidth}}{
        \cellcolor[rgb]{0.9,0.9,0.9}{
            \makecell[{{p{\linewidth}}}]{
                \texttt{\tiny{[GM$|$GM]}}
                \texttt{P moved to (7, 7)} \\
            }
        }
    }
    & & \\ \\

    \theutterance \stepcounter{utterance}  
    & & \multicolumn{4}{p{0.6\linewidth}}{
        \cellcolor[rgb]{0.9,0.9,0.9}{
            \makecell[{{p{\linewidth}}}]{
                \texttt{\tiny{[P1$\langle$GM]}}
                \texttt{Déplacement de 'O' vers (7,7) réussi. Votre grille mise à jour ressemble à ceci :} \\
\texttt{```} \\
\texttt{ 1234567} \\
\texttt{╔══╤══╤═╗} \\
\texttt{║◌◌│◌◌│P║ 1} \\
\texttt{╟──┼──┴─╢ 2} \\
\texttt{║◌◌│◌ST◌║ 3} \\
\texttt{╟──┘◌◌◌◌║ 4} \\
\texttt{║◌◌◌◌◌DA║ 5} \\
\texttt{║◌◌┌─┬──╢ 6} \\
\texttt{║◌◌│◌│MO║ 7} \\
\texttt{╚══╧═╧══╝} \\
\texttt{```} \\
\texttt{Vous jouez actuellement le tour 10 sur un maximum de 28.} \\
\texttt{Vous avez collectivement accumulé 21 pénalités sur 16.} \\
\texttt{L'autre joueur a déplacé un objet sur sa grille.} \\
\\ 
\\ 
\texttt{Quelle est votre prochaine commande ?} \\
            }
        }
    }
    & \\ \\

    \theutterance \stepcounter{utterance}  
    & \multicolumn{4}{p{0.6\linewidth}}{
        \cellcolor[rgb]{0.9,0.9,0.9}{
            \makecell[{{p{\linewidth}}}]{
                \texttt{\tiny{[P1$\rangle$GM]}}
                \texttt{DÉPLACER: P, (6, 7)} \\
            }
        }
    }
    & & \\ \\

    \theutterance \stepcounter{utterance}  
    & & & \multicolumn{2}{p{0.3\linewidth}}{
        \cellcolor[rgb]{0.9,0.9,0.9}{
            \makecell[{{p{\linewidth}}}]{
                \texttt{\tiny{[GM$|$GM]}}
                \texttt{DÉPLACER: P, (6, 7)} \\
            }
        }
    }
    & & \\ \\

    \theutterance \stepcounter{utterance}  
    & & & \multicolumn{2}{p{0.3\linewidth}}{
        \cellcolor[rgb]{0.9,0.9,0.9}{
            \makecell[{{p{\linewidth}}}]{
                \texttt{\tiny{[GM$|$GM]}}
                \texttt{Pénalité : (6,7) n'est pas vide, mais contient 'M'.} \\
\\ 
\texttt{Vous avez collectivement accumulé 22 pénalités sur 16.} \\
\texttt{Veuillez réessayer !} \\
            }
        }
    }
    & & \\ \\

    \theutterance \stepcounter{utterance}  
    & & \multicolumn{4}{p{0.6\linewidth}}{
        \cellcolor[rgb]{0.9,0.9,0.9}{
            \makecell[{{p{\linewidth}}}]{
                \texttt{\tiny{[P1$\langle$GM]}}
                \texttt{Pénalité : (6,7) n'est pas vide, mais contient 'M'.} \\
\\ 
\texttt{Vous avez collectivement accumulé 22 pénalités sur 16.} \\
\texttt{Veuillez réessayer !} \\
            }
        }
    }
    & \\ \\

    \theutterance \stepcounter{utterance}  
    & \multicolumn{4}{p{0.6\linewidth}}{
        \cellcolor[rgb]{0.9,0.9,0.9}{
            \makecell[{{p{\linewidth}}}]{
                \texttt{\tiny{[P1$\rangle$GM]}}
                \texttt{DÉPLACER: P, (6, 5)} \\
            }
        }
    }
    & & \\ \\

    \theutterance \stepcounter{utterance}  
    & & & \multicolumn{2}{p{0.3\linewidth}}{
        \cellcolor[rgb]{0.9,0.9,0.9}{
            \makecell[{{p{\linewidth}}}]{
                \texttt{\tiny{[GM$|$GM]}}
                \texttt{DÉPLACER: P, (6, 5)} \\
            }
        }
    }
    & & \\ \\

    \theutterance \stepcounter{utterance}  
    & & & \multicolumn{2}{p{0.3\linewidth}}{
        \cellcolor[rgb]{0.9,0.9,0.9}{
            \makecell[{{p{\linewidth}}}]{
                \texttt{\tiny{[GM$|$GM]}}
                \texttt{Pénalité : (6,5) n'est pas vide, mais contient 'D'.} \\
\\ 
\texttt{Vous avez collectivement accumulé 23 pénalités sur 16.} \\
\texttt{Veuillez réessayer !} \\
            }
        }
    }
    & & \\ \\

    \theutterance \stepcounter{utterance}  
    & & \multicolumn{4}{p{0.6\linewidth}}{
        \cellcolor[rgb]{0.9,0.9,0.9}{
            \makecell[{{p{\linewidth}}}]{
                \texttt{\tiny{[P1$\langle$GM]}}
                \texttt{Pénalité : (6,5) n'est pas vide, mais contient 'D'.} \\
\\ 
\texttt{Vous avez collectivement accumulé 23 pénalités sur 16.} \\
\texttt{Veuillez réessayer !} \\
            }
        }
    }
    & \\ \\

    \theutterance \stepcounter{utterance}  
    & \multicolumn{4}{p{0.6\linewidth}}{
        \cellcolor[rgb]{0.9,0.9,0.9}{
            \makecell[{{p{\linewidth}}}]{
                \texttt{\tiny{[P1$\rangle$GM]}}
                \texttt{DÉPLACER: P, (7, 5)} \\
            }
        }
    }
    & & \\ \\

    \theutterance \stepcounter{utterance}  
    & & & \multicolumn{2}{p{0.3\linewidth}}{
        \cellcolor[rgb]{0.9,0.9,0.9}{
            \makecell[{{p{\linewidth}}}]{
                \texttt{\tiny{[GM$|$GM]}}
                \texttt{DÉPLACER: P, (7, 5)} \\
            }
        }
    }
    & & \\ \\

    \theutterance \stepcounter{utterance}  
    & & & \multicolumn{2}{p{0.3\linewidth}}{
        \cellcolor[rgb]{0.9,0.9,0.9}{
            \makecell[{{p{\linewidth}}}]{
                \texttt{\tiny{[GM$|$GM]}}
                \texttt{Pénalité : (7,5) n'est pas vide, mais contient 'A'.} \\
\\ 
\texttt{Vous avez collectivement accumulé 24 pénalités sur 16.} \\
\texttt{Veuillez réessayer !} \\
            }
        }
    }
    & & \\ \\

    \theutterance \stepcounter{utterance}  
    & & \multicolumn{4}{p{0.6\linewidth}}{
        \cellcolor[rgb]{0.9,0.9,0.9}{
            \makecell[{{p{\linewidth}}}]{
                \texttt{\tiny{[P1$\langle$GM]}}
                \texttt{Pénalité : (7,5) n'est pas vide, mais contient 'A'.} \\
\\ 
\texttt{Vous avez collectivement accumulé 24 pénalités sur 16.} \\
\texttt{Veuillez réessayer !} \\
            }
        }
    }
    & \\ \\

    \theutterance \stepcounter{utterance}  
    & \multicolumn{4}{p{0.6\linewidth}}{
        \cellcolor[rgb]{0.9,0.9,0.9}{
            \makecell[{{p{\linewidth}}}]{
                \texttt{\tiny{[P1$\rangle$GM]}}
                \texttt{DIRE: Je n'ai plus de cases vides dans le coin bas{-}droit (6{-}7, 5{-}7) : D(6,5), A(7,5), M(6,7), O(7,7). Il ne reste que (7,6) mais c'est un mur '─'. Je dois utiliser une autre zone. Je vais déplacer S (libre en (4,3)) vers (1,6) qui est vide, et T de (6,3) vers (1,7). Puis P viendra en (6,3) libéré. Nouvelle cible : zone gauche lignes 6{-}7. On s'aligne là{-}bas ?} \\
            }
        }
    }
    & & \\ \\

    \theutterance \stepcounter{utterance}  
    & & & \multicolumn{2}{p{0.3\linewidth}}{
        \cellcolor[rgb]{0.9,0.9,0.9}{
            \makecell[{{p{\linewidth}}}]{
                \texttt{\tiny{[GM$|$GM]}}
                \texttt{DIRE: Je n'ai plus de cases vides dans le coin bas{-}droit (6{-}7, 5{-}7) : D(6,5), A(7,5), M(6,7), O(7,7). Il ne reste que (7,6) mais c'est un mur '─'. Je dois utiliser une autre zone. Je vais déplacer S (libre en (4,3)) vers (1,6) qui est vide, et T de (6,3) vers (1,7). Puis P viendra en (6,3) libéré. Nouvelle cible : zone gauche lignes 6{-}7. On s'aligne là{-}bas ?} \\
            }
        }
    }
    & & \\ \\

    \theutterance \stepcounter{utterance}  
    & & \multicolumn{4}{p{0.6\linewidth}}{
        \cellcolor[rgb]{0.9,0.9,0.9}{
            \makecell[{{p{\linewidth}}}]{
                \texttt{\tiny{[GM$\rangle$P2]}}
                \texttt{Déplacement de 'P' vers (7,7) réussi. Votre grille mise à jour ressemble à ceci :} \\
\texttt{```} \\
\texttt{ 1234567} \\
\texttt{╔══╤══╤═╗} \\
\texttt{║◌◌│◌◌│◌║ 1} \\
\texttt{╟──┼──┴─╢ 2} \\
\texttt{║◌◌│◌◌◌◌║ 3} \\
\texttt{╟──┘◌◌T◌║ 4} \\
\texttt{║D◌◌◌◌MO║ 5} \\
\texttt{║◌◌┌─┬──╢ 6} \\
\texttt{║◌◌│A│SP║ 7} \\
\texttt{╚══╧═╧══╝} \\
\texttt{```} \\
\texttt{Vous jouez actuellement le tour 11 sur un maximum de 28.} \\
\texttt{Vous avez collectivement accumulé 24 pénalités sur 16.} \\
\texttt{L'autre joueur a envoyé ce message :} \\
\\ 
\\ 
\texttt{"Je n'ai plus de cases vides dans le coin bas{-}droit (6{-}7, 5{-}7) : D(6,5), A(7,5), M(6,7), O(7,7). Il ne reste que (7,6) mais c'est un mur '─'. Je dois utiliser une autre zone. Je vais déplacer S (libre en (4,3)) vers (1,6) qui est vide, et T de (6,3) vers (1,7). Puis P viendra en (6,3) libéré. Nouvelle cible : zone gauche lignes 6{-}7. On s'aligne là{-}bas ?"} \\
\\ 
\\ 
\texttt{Quelle est votre prochaine commande ?} \\
            }
        }
    }
    & \\ \\

    \theutterance \stepcounter{utterance}  
    & & & \multicolumn{4}{p{0.6\linewidth}}{
        \cellcolor[rgb]{0.9,0.9,0.9}{
            \makecell[{{p{\linewidth}}}]{
                \texttt{\tiny{[GM$\langle$P2]}}
                \texttt{DIRE: D'accord pour la zone gauche lignes 6{-}7. Chez moi : D(1,5), A(3,7), T(6,4), M(6,5), O(7,5), S(6,7), P(7,7). Cases libres col 1{-}2 lig 6{-}7 : (1,6), (2,6), (2,7). (1,7) est mur. On a 3 cases pour 7 objets, c'est trop peu. Il faut plus de place. } \\
\\ 
\texttt{Quelles cases vides as{-}tu lignes 5{-}7, col 1{-}4 ? Moi j'ai (2,5),(3,5),(4,5),(5,5),(6,5) prise,(7,5) prise, (1,6),(2,6),(3,6),(4,6),(5,6),(6,6) mur,(7,6) mur, (2,7). } \\
\\ 
\texttt{On peut viser un rectangle 2x4 ou 3x3 quelque part. Propose une zone précise avec coordonnées.} \\
            }
        }
    }
     \\ \\

    \theutterance \stepcounter{utterance}  
    & & & \multicolumn{2}{p{0.3\linewidth}}{
        \cellcolor[rgb]{0.9,0.9,0.9}{
            \makecell[{{p{\linewidth}}}]{
                \texttt{\tiny{[GM$|$GM]}}
                \texttt{DIRE: D'accord pour la zone gauche lignes 6{-}7. Chez moi : D(1,5), A(3,7), T(6,4), M(6,5), O(7,5), S(6,7), P(7,7). Cases libres col 1{-}2 lig 6{-}7 : (1,6), (2,6), (2,7). (1,7) est mur. On a 3 cases pour 7 objets, c'est trop peu. Il faut plus de place. } \\
\\ 
\texttt{Quelles cases vides as{-}tu lignes 5{-}7, col 1{-}4 ? Moi j'ai (2,5),(3,5),(4,5),(5,5),(6,5) prise,(7,5) prise, (1,6),(2,6),(3,6),(4,6),(5,6),(6,6) mur,(7,6) mur, (2,7). } \\
\\ 
\texttt{On peut viser un rectangle 2x4 ou 3x3 quelque part. Propose une zone précise avec coordonnées.} \\
            }
        }
    }
    & & \\ \\

    \theutterance \stepcounter{utterance}  
    & & & \multicolumn{2}{p{0.3\linewidth}}{
        \cellcolor[rgb]{0.9,0.9,0.9}{
            \makecell[{{p{\linewidth}}}]{
                \texttt{\tiny{[GM$|$GM]}}
                \texttt{Invalid format: tail is not empty: '} \\
\\ 
\texttt{Quelles cases vides as{-}tu lignes 5{-}7, col 1{-}4 ? Moi j'ai (2,5),(3,5),(4,5),(5,5),(6,5) prise,(7,5) prise, (1,6),(2,6),(3,6),(4,6),(5,6),(6,6) mur,(7,6) mur, (2,7). } \\
\\ 
\texttt{On peut viser un rectangle 2x4 ou 3x3 quelque part. Propose une zone précise avec coordonnées.} \\
            }
        }
    }
    & & \\ \\

    \theutterance \stepcounter{utterance}  
    & & & \multicolumn{2}{p{0.3\linewidth}}{
        \cellcolor[rgb]{0.9,0.9,0.9}{
            \makecell[{{p{\linewidth}}}]{
                \texttt{\tiny{[GM$|$GM]}}
                \texttt{Maximum number of penalties exceeded} \\
            }
        }
    }
    & & \\ \\

    \theutterance \stepcounter{utterance}  
    & & & \multicolumn{2}{p{0.3\linewidth}}{
        \cellcolor[rgb]{0.9,0.9,0.9}{
            \makecell[{{p{\linewidth}}}]{
                \texttt{\tiny{[GM$|$GM]}}
                \texttt{Initial states:} \\
\texttt{```} \\
\texttt{Player 1:} \\
\texttt{ 1234567} \\
\texttt{╔══╤══╤═╗} \\
\texttt{║◌P│◌M│◌║ 1} \\
\texttt{╟──┼──┴─╢ 2} \\
\texttt{║◌◌│◌S◌O║ 3} \\
\texttt{╟──┘◌A◌◌║ 4} \\
\texttt{║D◌◌◌◌◌◌║ 5} \\
\texttt{║◌◌┌─┬──╢ 6} \\
\texttt{║T◌│◌│◌◌║ 7} \\
\texttt{╚══╧═╧══╝} \\
\\ 
\\ 
\texttt{Player 2:} \\
\texttt{ 1234567} \\
\texttt{╔══╤══╤═╗} \\
\texttt{║◌◌│◌◌│T║ 1} \\
\texttt{╟──┼──┴─╢ 2} \\
\texttt{║◌◌│◌S◌P║ 3} \\
\texttt{╟──┘◌◌MO║ 4} \\
\texttt{║D◌◌◌◌◌◌║ 5} \\
\texttt{║◌◌┌─┬──╢ 6} \\
\texttt{║◌◌│A│◌◌║ 7} \\
\texttt{╚══╧═╧══╝} \\
\\ 
\texttt{```} \\
            }
        }
    }
    & & \\ \\

    \theutterance \stepcounter{utterance}  
    & & & \multicolumn{2}{p{0.3\linewidth}}{
        \cellcolor[rgb]{0.9,0.9,0.9}{
            \makecell[{{p{\linewidth}}}]{
                \texttt{\tiny{[GM$|$GM]}}
                \texttt{End states:} \\
\texttt{```} \\
\texttt{Player 1:} \\
\texttt{ 1234567} \\
\texttt{╔══╤══╤═╗} \\
\texttt{║◌◌│◌◌│P║ 1} \\
\texttt{╟──┼──┴─╢ 2} \\
\texttt{║◌◌│◌ST◌║ 3} \\
\texttt{╟──┘◌◌◌◌║ 4} \\
\texttt{║◌◌◌◌◌DA║ 5} \\
\texttt{║◌◌┌─┬──╢ 6} \\
\texttt{║◌◌│◌│MO║ 7} \\
\texttt{╚══╧═╧══╝} \\
\\ 
\\ 
\texttt{Player 2:} \\
\texttt{ 1234567} \\
\texttt{╔══╤══╤═╗} \\
\texttt{║◌◌│◌◌│◌║ 1} \\
\texttt{╟──┼──┴─╢ 2} \\
\texttt{║◌◌│◌◌◌◌║ 3} \\
\texttt{╟──┘◌◌T◌║ 4} \\
\texttt{║D◌◌◌◌MO║ 5} \\
\texttt{║◌◌┌─┬──╢ 6} \\
\texttt{║◌◌│A│SP║ 7} \\
\texttt{╚══╧═╧══╝} \\
\\ 
\texttt{```} \\
            }
        }
    }
    & & \\ \\

    \theutterance \stepcounter{utterance}  
    & & & \multicolumn{2}{p{0.3\linewidth}}{
        \cellcolor[rgb]{0.9,0.9,0.9}{
            \makecell[{{p{\linewidth}}}]{
                \texttt{\tiny{[GM$|$GM]}}
                \texttt{* success: False} \\
\texttt{* lose: False} \\
\texttt{* aborted: True} \\
\texttt{{-}{-}{-}{-}{-}{-}{-}} \\
\texttt{* Shifts: 10.00} \\
\texttt{* Max Shifts: 12.00} \\
\texttt{* Min Shifts: 6.00} \\
\texttt{* End Distance Sum: 23.73} \\
\texttt{* Init Distance Sum: 21.20} \\
\texttt{* Expected Distance Sum: 29.33} \\
\texttt{* Penalties: 25.00} \\
\texttt{* Max Penalties: 16.00} \\
\texttt{* Rounds: 10.00} \\
\texttt{* Max Rounds: 28.00} \\
\texttt{* Object Count: 7.00} \\
            }
        }
    }
    & & \\ \\

\end{supertabular}
}

\end{document}
